\ifdefined\ICLRSubmission
  \documentclass{article}
  \usepackage{iclr2027/iclr2027_conference,times}
\else
  \documentclass[10pt]{article}
  \usepackage[margin=1in]{geometry}
\fi

\usepackage{amsmath,amssymb,amsthm,mathtools}
\usepackage{booktabs,multirow}
\usepackage{graphicx}
\usepackage[dvipsnames]{xcolor}
\usepackage{microtype}
\usepackage{enumitem}
\usepackage{natbib}
\usepackage[hidelinks]{hyperref}
\usepackage{url}

\ifdefined\ICLRSubmission
  % Keep the official under-review banner visible with current fancyhdr releases.
  % The vendored conference style itself remains byte-identical to the official file.
  \lhead{Under review as a conference paper at ICLR 2027}
\fi

\newtheorem{theorem}{Theorem}
\newtheorem{corollary}[theorem]{Corollary}
\newtheorem{proposition}[theorem]{Proposition}
\newtheorem{lemma}[theorem]{Lemma}
\theoremstyle{definition}
\newtheorem{definition}[theorem]{Definition}
\newtheorem{assumption}[theorem]{Assumption}
\theoremstyle{remark}
\newtheorem{remark}[theorem]{Remark}

\newcommand{\R}{\mathbb{R}}
\newcommand{\E}{\mathbb{E}}
\newcommand{\Prob}{\mathbb{P}}
\newcommand{\ip}[2]{\left\langle #1,#2\right\rangle}
\newcommand{\norm}[1]{\left\lVert #1\right\rVert}
\newcommand{\given}{\,\middle|\,}
\newcommand{\sigm}{\sigma}
\newcommand{\ind}{\mathbf{1}}
\newcommand{\Bcond}{\mathrm{B}}
\newcommand{\Wcond}{\mathrm{W}}
\newcommand{\dd}{\,\mathrm{d}}
\newcommand{\code}[1]{\texttt{#1}}

\newcommand{\artifactfigure}[3]{%
\begin{figure}[t]
  \centering
  \IfFileExists{#1}{\includegraphics[width=0.97\linewidth]{#1}}{%
    \fbox{\parbox{0.92\linewidth}{\centering Figure artifact not found while
    compiling this source. Expected path: \texttt{\detokenize{#1}}.}}}
  \caption{#2}
  \label{#3}
\end{figure}}

\title{Weak Learning Failure Alone Is Not Gradient Starvation:\\
Exact Finite-Width Dynamics for a Paired Causal Test}
\ifdefined\ICLRSubmission
  \author{}
\else
  \author{Anonymous authors}
  \date{Working draft}
\fi

\begin{document}
\maketitle

\begin{abstract}
A weak feature can remain poorly learned for several different reasons: its
weak-only trajectory may not reach a specified target within the observation
horizon, a strong cue may change cross-entropy weighting, response-specific
parameter geometry may suppress its update, or earlier optimization help may not
yet have been repaid by later suppression.  Weak learning failure alone is not
gradient starvation.  We therefore compare a both-feature trajectory with a
same-initialization weak-only trajectory in which only the strong input channel is
ablated.  This paired intervention distinguishes relative transfer, rate
suppression, outcome suppression, weak-only nonlearnability under the specified
target and horizon, and causal starvation.  A failed weak-only first-hit gate
licenses neither a starvation certificate nor evidence that starvation is absent;
the causal status is indeterminate under that gate.

For full-batch logistic gradient flow, the exact finite-width weak-response
dynamics are $\dot M_w=c_w^\top q(r)/n$, where $q_i(r)=\sigma(-r_i)$ is the
cross-entropy weighting of signed logit $r_i$ and $c_w=J\nabla M_w$ is the
response-specific response--logit geometry.  The companion signed-logit flow is
$\dot r=\Theta q(r)/n$.  These chain-rule identities require no dynamical
mean-field theory (DMFT) or response-coordinate closure.  Freezing
$(\Theta,c_w)$ defines a nonlinear logistic tangent surrogate; continuous
Gr\"onwall and experiment-matched discrete secant-kernel bounds compare it with
nonlinear training only under supplied kernel-movement control.  Pathwise, for
arbitrary sign changes, $\Delta=P-N$: outcome suppression occurs exactly when
accumulated negative relative drift overtakes accumulated positive relative drift.
Causal starvation additionally requires independently demonstrated weak-only
learnability.  A noiseless rank-one log drift ratio gives sufficient conditions for
one rate crossing.  In the dense-linear model, the realized initialization gives
an exact pre-trajectory $d(0)$ threshold and an existential strict local outcome
sign outside a caller indifference band; where proved, local jet information such
as $d'(0)$ supplies further assumption-scoped local signs.  Uniform product-ball
bounds, when independently proved, imply crossing, outcome, nondegenerate
learnability, or finite-horizon safety, but no such nonvacuous general positive-lag
constants are currently constructed.  Thus stronger global parameter-level
prediction from $\rho$, lag, and recurrent geometry remains unresolved.  We also
prove stability of quantitatively isolated transverse first hits under uniform-
error, radius, and mesh hypotheses and an exact six-scalar zero-lag wide limit for
any finite realized recurrent matrix whose noiseless symmetric reference is
suppressed from initialization while its weak-only trajectory remains learnable.
As local mechanism-level control, in the fixed implemented chart, feasible
uncapped target-attaining Counterfactual Drift Correction (CDC) is the unique
minimum-Euclidean-norm correction satisfying zero instantaneous strong-drift
change and the weak-drift lower bound; its finite-step guarantee is local and
one-step.

In controlled dense-linear E1, all 32 $\rho$--seed cells per lag at lags 0, 2,
and 4 (four strengths by eight seed IDs) have positive paired weak-AUC deficits and
learnable weak-only controls; at each fixed lag, the mean deficit increases over
the measured $\rho$ grid.  All 32 lag-8 cells fail the weak-only gate and are
causally indeterminate under $\beta=0.25,H=5$.  At the fixed
$(\rho,\mathrm{lag})=(4,2)$ nonlinear cell, tanh has complete transfer-to-starvation
certificates in 3/8 observed seeds; GRU has drift and response crossings in 5/8 but
0/8 causal certificates because every weak-only path fails $\beta=0.5$.  A
post-hoc terminal sensitivity audit, derived without retraining, places tanh
$B_W(H)$ in $[4.413,4.676]$ and GRU in $[0.0828,0.187]$; it is not a replacement
for the first-hit gate or an architecture ranking.  A sealed five-endpoint frozen-
kernel E19 pilot fails its preregistered rule.  Only after that failure, E20's fresh
restricted $4\times8$ data/model-seed factorial classifies response-crossing events
in 64/64 records and drift-crossing events in 60/64, with two-way-bootstrap 95\%
lower endpoints 1.0 and 0.8125.  These are one-cell event-classification results,
not phase, causal, learnability, calibrated-time, magnitude, prevalence, or kernel-
stability validation.  E20 remains the submission-primary empirical study with
only its frozen drift/response-crossing endpoints primary; later clean-source
extensions do not reprioritize it.  E26 has 64 semi-real generated-cue records:
MNIST 3-vs-8 and FashionMNIST 0-vs-6 each have 0/32 weak-only-learnable and 0/32
causal records, with mean weak-AUC gaps
$-1.2362842136667493\times10^{-3}$ (95\% CI
$[-3.4405428466004646,1.0159591371180453]\times10^{-3}$) and
$-1.0921728159018996\times10^{-4}$ (95\% CI
$[-4.129384520886106,3.182670049955049]\times10^{-3}$), respectively; this is
negative/indeterminate.  The failed causal gate blocks a false-positive starvation
interpretation but licenses no absence claim.  In E27's 192 records, tanh's
all-$\beta$ estimate is $0.2916666666666667$ (95\% CI
$[0.16666666666666666,0.4166666666666667]$) and GRU's is $0$ (95\% CI
$[0,0]$), with Holm-adjusted $p=1$ for each; the frozen endpoint is negative and
supports no architecture ranking.  E28's 256 method-records pass the frozen joint
weak-rescue/strong-preservation tradeoff only against the non-canonical Bloop-style
shadow-target rescue; CDC does not beat the non-canonical PCGrad-style shadow-
target rescue under the frozen joint criterion.  No SOTA, canonical parity, or
broad CDC superiority follows.  The resulting chain is exact dynamics $\to$
matched causal intervention $\to$ pathwise attribution $\to$ local mechanism-level
control.  A general positive-lag optimization-time DMFT remains open.
\end{abstract}

\section{Introduction}
A weak predictive feature can remain small for several distinct reasons.  The
weak-only task may not reach a chosen target within the observation horizon; a
strong cue may reduce cross-entropy weight on already-correct samples; the
response-specific parameter geometry may route little gradient toward the weak
response; or a later period of rate suppression may be too short to erase earlier
help.  Cross-entropy can indeed favor one predictive feature strongly enough to
starve another \citep{pezeshki2021gradient}, but weak learning failure alone is not
gradient starvation.  Recurrent networks make the diagnosis especially delicate
because cue strength, lag, shared recurrent parameters, and evolving response
geometry all affect the update.

We make the diagnosis with a matched intervention.  The both-feature and weak-only
models share initialization, labels, samples, weak coordinates, and noise; only the
strong input channel is set to zero in the weak-only condition.  We track common
feature-response functionals in both trajectories.  This setup separates five
statements:
\begin{enumerate}[leftmargin=*,nosep]
  \item \emph{relative transfer}: the strong cue initially increases the weak
  response rate, so $d(\tau)>0$;
  \item \emph{rate suppression}: the both-feature weak response later grows more
  slowly than its counterfactual, so $d(\tau)<0$;
  \item \emph{outcome suppression}: the both-feature weak response falls below
  its matched counterfactual, so $\Delta(\tau)<0$;
  \item \emph{weak-only nonlearnability under the specified gate}: the weak-only
  path does not first hit the preregistered target by the horizon; and
  \item \emph{causal starvation}: outcome suppression occurs while the weak-only
  counterfactual demonstrably learns under that first-hit gate.
\end{enumerate}
A failed weak-only gate therefore licenses neither a causal-starvation certificate
nor evidence that starvation is absent; causal status is indeterminate under the
specified target and horizon.  A transfer-to-starvation trajectory is a particular
causally certified path with an initial positive excursion before outcome
suppression, not a default description of weak learning failure.  This distinction
is empirical as well as logical: at the fixed nonlinear cell, tanh has complete
transfer-to-starvation certificates in 3/8 observed seeds, whereas GRU has drift
and response crossings in 5/8 but fails the causal learnability gate in every seed.
These counts do not rank architectures or estimate prevalence.

The diagnostic has an exact finite-width dynamical basis.  For signed sample logits
$r$, full-batch logistic gradient flow gives
$\dot M_w=c_w^\top q(r)/n$, with $q_i(r)=\sigma(-r_i)$ the cross-entropy weighting
and $c_w=J\nabla M_w$ the response-specific response--logit geometry.  No DMFT is
needed for this chain-rule identity.  For the paired gap, the pathwise balance
$\Delta=P-N$ allows arbitrary drift sign changes: outcome suppression occurs only
when accumulated suppression overtakes accumulated help, and causal starvation
still requires the independent weak-only gate.  Prediction is partial.  The
realized-initialization value $d(0)$ and, where proved, local jet information such
as $d'(0)$ give rigorous local or assumption-scoped sufficient signs.  Stronger
global parameter-level prediction from $\rho$, lag, and RNN geometry remains
unresolved because nonvacuous positive-lag product-ball and kernel-movement
constants are unavailable.  CDC then uses the same matched weak-only information
as a local mechanism-level control, not as a claimed general remedy.

\paragraph{Contributions.}
\begin{enumerate}[leftmargin=*]
  \item We operationalize a paired causal test that separates relative transfer,
  rate suppression, outcome suppression, weak-only gate failure, and causal
  starvation through direct drift, first-hit, AUC, and finite-step certificates.
  A post-hoc terminal $\beta$ profile reports threshold sensitivity from existing
  outcomes without replacing the preregistered first-hit gate.
  \item We report the full findings sequence rather than only positive endpoints.
  Controlled E1 has positive paired weak-AUC deficits and learnable weak-only
  controls in all measured cells at lags 0, 2, and 4, while lag 8 fails the gate
  and is causally indeterminate.  E-NL separates 3/8 tanh transfer-to-starvation
  certificates from 5/8 ungated GRU drift/response crossings.  E19 remains a
  failed five-endpoint pilot; E20 remains the submission-primary empirical claim
  and supports only fresh restricted drift/response event labels.  E26 is
  negative/indeterminate after all records fail their specified gate; E27 is
  negative without an architecture ranking; and E28 yields a comparator-restricted
  mixed CDC tradeoff.  Archives distinguish historically unreconstructible source
  from clean, reconstructible source and from self-contained replay packages.
  \item We derive exact finite-width logistic response flow and full signed-logit
  empirical-NTK/response dynamics.  The weak response obeys
  $\dot M_w=c_w^\top q(r)/n$, explicitly separating cross-entropy weighting from
  response-specific geometry.  A closed projected equation $\dot M=Gg$ requires
  exact parameter-independent logit realization; the initialization-frozen
  nonlinear logistic surrogate is exact internally, while comparison with a
  trained nonlinear model is conditional on supplied kernel-movement bounds.
  \item We prove the pathwise cumulative drift balance for arbitrary sign changes,
  with the exact single-crossing tail-area criterion as a corollary.  We add a
  conditional noiseless rank-one rate-crossing criterion and an exact realized-
  initialization dense-linear $d(0)$ threshold with its existential strict local
  outcome-sign implication.  Independently proved product-ball bounds then imply
  unique crossing, outcome suppression, nondegenerate weak-only learnability, or
  finite-horizon safe transfer; no current component constructs these bounds for
  general positive-lag RNNs.  We also prove stability of quantitatively isolated
  transverse first hits under uniform-error, radius, and mesh hypotheses, and an
  exact zero-lag recurrence-invisible six-scalar wide limit for arbitrary finite
  realized recurrent gain, with immediate deterministic suppression and weak-only
  learnability under the noiseless cue law.
  \item As local mechanism-level control, in the fixed implemented
  Euclidean/Frobenius chart, feasible uncapped target-attaining CDC is the unique
  minimum-Euclidean-norm correction satisfying zero instantaneous strong-drift
  change and the weak-drift lower bound.  Its finite-step guarantee is local and
  one-step.  The ablations establish neither SOTA performance, canonical parity,
  broad practical superiority, nor best final retention.
\end{enumerate}

\paragraph{Claim boundary.}
We do not claim that recurrent networks universally undergo transfer-to-starvation,
that a general joint CE-RNN optimization-time DMFT has been solved, or that an
analytic $\rho_c$ or universal $(\rho,\mathrm{lag})$ phase boundary is known.  A
failed weak-only gate is causal indeterminacy under the specified target and
horizon: it establishes neither causal starvation nor its absence, and it does not
show that a dataset is intrinsically unlearnable.  The exact initialization
threshold and proved local jets provide local or assumption-scoped signs, not a
general predictor; at positive lag, the threshold is realization-dependent and its
trajectory implications require uniform product-ball constants that are not
currently constructed.  The frozen empirical-NTK surrogate is exact for its
initialization-linearized tangent model, not for trained tanh/GRU networks: no
required instantaneous or secant-kernel stability bound is proved here.  E19 is a
failed broad pilot.  E20 validates only binary drift/response-crossing
classification at one width/task cell; it does not validate phase, causal
starvation, weak-only learnability, calibrated trajectories, crossing times,
magnitudes, or prevalence.  E20 remains submission-primary with only its frozen
drift/response-crossing endpoints primary; E26--E28 do not outcome-dependently
reprioritize it.  The eight-seed nonlinear proportions are descriptive, not
population prevalence.  E26 uses real benchmark pixels with a generated removable
cue and therefore supplies semi-real, not natural-spurious-correlation, evidence;
its absent weak-only first hits make zero causal certificates indeterminate.  E27
supports no beta-robust architecture claim or architecture ranking.  E28 supports
a joint tradeoff only against one of two non-canonical comparators; CDC does not
beat the non-canonical PCGrad-style shadow-target rescue under the frozen joint
criterion.  No SOTA, broad CDC superiority, or canonical parity is claimed.  The
general DMFT solver remains deliberately unimplemented where proof obligations are
unmet.

\section{Related work}
\paragraph{Gradient starvation.}
\citet{pezeshki2021gradient} identify feature starvation under cross-entropy and
motivate Spectral Decoupling.  Our focus is a recurrent paired causal diagnosis:
the same-initialization weak-only intervention tests the learnability premise, exact
finite-width dynamics separate cross-entropy weighting from response-specific
geometry, and pathwise balance separates rate suppression from outcome suppression
and causally gated starvation.  A failed weak-only gate is indeterminate rather
than evidence for either starvation or its absence.

\paragraph{Shortcut learning and spurious-correlation robustness.}
Shortcut learning and simplicity bias explain why predictive but unintended cues
can dominate learned representations \citep{geirhos2020shortcut,shah2020pitfalls}.
Group DRO, Just Train Twice, and last-layer retraining instead improve worst-group
performance under spurious correlations
\citep{sagawa2020groupdro,liu2021jtt,kirichenko2023dfr}.  These phenomena motivate
our setting, but they do not supply the paired response-level causal criterion used
here.  We are not proposing a competing robustness algorithm: our diagnostic
requires a stronger matched-shadow information setting.  E26 adds a generated-cue
intervention on real MNIST and FashionMNIST core pixels, but it is not an unmodified
natural-spurious-correlation benchmark and supports no such claim.

\paragraph{Learning dynamics in recurrent and wide networks.}
Neural tangent kernels describe infinite-width gradient dynamics
\citep{jacot2018ntk}, and wide-network linearization motivates frozen-tangent
comparisons \citep{lee2019wide}.  Our initialization-frozen logistic surrogate
retains the nonlinear loss but is exact only for its tangent model.  Its comparison
with a trained finite-width network is conditional on explicit instantaneous or
secant-kernel movement bounds, which we do not prove for tanh or GRU.
\citet{ger2026invisible} derive low-dimensional overlap learning dynamics for
low-rank RNNs and show that overlaps invisible to the loss can retain training
history.  This directly cautions against closing a recurrent learning theory using
only visible response coordinates.  \citet{clark2026structure} develop a
mean-field account of structure and disorder in task-trained recurrent circuits in
a long-time Langevin/Gibbs setting.  That equilibrium construction is distinct from
the paired, noiseless optimization-time gradient flow studied here.  Accordingly,
we claim neither the first learning theory for RNNs nor a completed general RNN
DMFT; our contribution is the matched causal separation supported by exact finite-
width response and pathwise attribution results.\footnote{The descriptions of the 2026 sources were paraphrased from the
linked public sources. Content was rephrased for compliance with licensing
restrictions.}

\paragraph{Gradient surgery.}
PCGrad projects conflicting task gradients \citep{yu2020pcgrad}; Bloop stabilizes
an auxiliary projection using an exponential moving average
\citep{hsieh2024bloop}.  Projection itself is not our novelty.  CDC protects the
gradient of a theory-defined strong-feature response and takes its target from a
matched causal weak-only shadow.  This stronger information setting must be
reported: CDC requires a constructible ablation, a second model, lockstep training,
and full-batch gradients.  In E28 the evaluated Bloop-style shadow-target rescue and
PCGrad-style shadow-target rescue are non-canonical (\code{canonical:false}); the
frozen joint weak-rescue/strong-preservation tradeoff holds only against the former,
and CDC does not beat the latter under that criterion.  The study establishes
neither SOTA, canonical parity, nor broad CDC superiority.

\ifdefined\ICLRSubmission
\section{Paired causal diagnosis}
Weak learning failure alone is not gradient starvation.  For condition
$C\in\{\Bcond,\Wcond\}$, the both-feature and weak-only models share
initialization, labels, samples, weak coordinates, and noise; only the strong input
channel is ablated under $\Wcond$.  The matched weak-only path tests whether the
weak signal can meet the specified target without the strong cue; it is an
identification premise rather than a second observation of the same failure.  Let
$M_w^C(\tau)$ be the common weak-feature response, and define
\begin{equation}
 \Delta(\tau)=M_w^{\Bcond}(\tau)-M_w^{\Wcond}(\tau),\qquad
 d(\tau)=\Delta'(\tau),\qquad \Delta(0)=0.
 \label{eq:iclr-gap}
\end{equation}
This pairing separates claims that are often conflated: $d>0$ is
\emph{relative transfer}; $d<0$ is \emph{rate suppression}; $\Delta<0$ is
\emph{outcome suppression}; failure to first hit a preregistered target $\beta$ by
horizon $H$ is \emph{weak-only nonlearnability under that gate}; and
\emph{causal starvation} requires outcome suppression together with a successful,
nondegenerate weak-only first hit.  A response crossing without weak-only
learnability is descriptive, not a causal certificate.  Conversely, the failed
gate is not evidence that starvation is absent; causal status is indeterminate
under the specified target and horizon.

For signed logits $r_i^C=y_if_C(x_i;\theta_C)$ and full-batch logistic gradient
flow, the exact finite-width chain rule is
\begin{equation}
 \dot M_a^C=\frac1{n_C}\sum_i
 \langle\nabla M_a^C,\nabla r_i^C\rangle\,\sigma(-r_i^C).
 \label{eq:iclr-cross-kernel}
\end{equation}
We state this as a proposition rather than a theorem: the derivation is elementary;
the contribution is its exact finite-width scope and the causal separation it
enables.  With sample Jacobian $J$, empirical NTK $\Theta=JJ^\top$, response
cross-kernel $c_a=J\nabla M_a$, and cross-entropy weights
$q_i(r)=\sigma(-r_i)$, the exact dynamics are
$\dot r=\Theta q(r)/n$ and $\dot M_a=c_a^\top q(r)/n$.  Thus weak-response
motion separates sample weighting $q(r)$ from response-specific geometry $c_w$;
no DMFT or response-coordinate closure is needed for this identity.  Freezing
$(\Theta,c_a)$ at initialization defines a nonlinear logistic tangent surrogate.
Its comparison with nonlinear training is conditional on supplied kernel-movement
bounds; event agreement is not a substitute for that premise.

\section{Pathwise separation of help, suppression, and causal status}
For any absolutely continuous $\Delta$, let
$P(t)=\int_0^t[d(u)]_+\,du$ and $N(t)=\int_0^t[-d(u)]_+\,du$.  Our general
bookkeeping proposition gives the exact identity
\begin{equation}
 \Delta(t)=P(t)-N(t),\qquad
 \Delta(t)<0\ \Longleftrightarrow\ N(t)>P(t).
 \label{eq:iclr-drift-balance}
\end{equation}
It allows arbitrary sign changes and immediate suppression: outcome suppression
occurs only when accumulated suppression $N$ overtakes accumulated help $P$.  Under
one $+\!\to\!-$ drift crossing at $\tau_d$, the identity reduces to the tail-area
criterion, in which negative post-crossing area must exceed the positive peak area.
Neither rate suppression nor outcome suppression establishes causal starvation
without the independent weak-only learnability gate; a failed gate leaves causal
status indeterminate.  The appendix proves this statement, a conditional rank-one
sufficient criterion for a unique rate crossing, and stability of quantitatively
isolated transverse first hits under uniform-error, radius, and mesh hypotheses.

The appendix also gives an initialization-conditioned pre-trajectory result.  In
the noiseless positive rank-one dense-linear pair, set
$K=G_{ws}(0)$, $A=G_{ww}(0)$,
$s_B=\sigma[-(\rho M_s(0)+M_w(0))]$, and
$s_W=\sigma[-M_w(0)]$.  The exact initial rate gap is
\begin{equation}
 d_0=s_B(\rho K+A)-s_WA,
 \quad d_0>0\Longleftrightarrow
 K>\frac{A}{\rho}\left(\frac{s_W}{s_B}-1\right).
 \label{eq:iclr-initial-gap}
\end{equation}
This is partial, initialization-conditioned prediction: it uses the realized
$d_0$, not a trained trajectory or a universal $(\rho,\mathrm{lag})$ boundary.
Shared initialization has $\Delta(0)=0$, so $d_0<0$ and $d_0>0$ respectively
imply strict local outcome suppression and transfer for all sufficiently small
positive times.  This derivative argument is existential and gives no quantitative
horizon; the executable certificate reports values within its caller tolerance as
boundary/undetermined.  Under independently supplied uniform product-ball bounds
$-\Lambda\le d'\le-\lambda<0$ up to $h_R=R/S_R$, one obtains
$d_0/\Lambda\le\tau_d\le d_0/\lambda$, and
$h_R>2d_0/\lambda$ implies outcome suppression.  A separate positive weak-only
drift bound gives causal learnability only for a target strictly above the initial
weak response; a target already met at initialization is degenerate and cannot
gate causal starvation.  Alternatively, $|d'|\le J$ and $d_0-Jh_R>0$ exclude
suppression on that finite horizon.  These are local or assumption-scoped
sufficient signs, not a general parameter-level predictor.  No current
implementation constructs nonvacuous general positive-lag constants from $\rho$,
lag, and recurrent geometry.

A proved local-jet example occurs in the deterministic finite-horizon wide limit
at zero lag, zero background input, and arbitrary finite realized recurrent gain:
all pre-final states vanish, so $W$ is loss-invisible.  For the explicitly noiseless
($\texttt{cue\_noise}=0$) positive rank-one law, its symmetric reference has
$\bar d(0)=0$ and $\bar d'(0)=-\rho^2/2<0$, so suppression starts immediately
rather than after transfer; the weak-only response reaches every finite positive
target.  This is a recurrence-invisible special case, not a predictor for general
positive lag.  A general positive-lag optimization-time DMFT remains open.

As local mechanism-level control, Counterfactual Drift Correction (CDC) uses the
weak-only shadow as a local target.  In the fixed implemented Euclidean/Frobenius
parameter chart, let $a=\nabla m_s$, $p=\nabla m_w$, $v=-\nabla L$, and let $q$
be the projection of $p$ onto $a^\perp$.  For positive deficit
$\delta=[F_w^{\Wcond}-p\cdot v]_+$, uncapped CDC takes
$v'=v+(\delta/\|q\|^2)q$ when feasible.  In that feasible uncapped
and target-attaining case, it is the unique minimum-Euclidean-norm correction
satisfying zero instantaneous strong-drift change and the weak-drift lower bound.
Its finite-step strong-response guarantee is local and one-step, of order
$O(\eta^2)$ under local smoothness.  These are fixed-chart, full-batch statements,
not trajectory, SOTA, canonical-parity, or practical-superiority claims.

\section{Experiments and audit protocol}
The evidence combines controlled synthetic studies with one semi-real generated-cue
intervention.  E1 uses a dense linear RNN with four strengths and seed IDs 4--11:
each lag therefore has 32 $\rho$--seed cells but only eight independent seed IDs.
E-NL evaluates one $(\rho,\mathrm{lag})=(4,2)$ cell with eight seeds each for tanh
and GRU, paired lockstep training, direct-autograd response drift, horizon $H$, and
preregistered first-hit target $\beta=0.5$.  Historical E3 compares CDC with ERM
and seven mitigation arms; the matched-shadow methods use stronger information and
roughly twice the model work.

The externally sealed clean-source extensions were executed at commit
\code{9ebbc617a9aaad5c3c19e9d5fd4fa08f2d3517d0}.  E26 crosses four data seeds with
eight model seeds on each of MNIST 3-vs-8 and FashionMNIST 0-vs-6, retaining real
core pixels and ablating only an exactly removable generated cue.  E27 contains 192
nonlinear records over three fixed cells, tanh/GRU, and prospective
$\beta\in\{0.25,0.5,1,2\}$.  E28 contains 256 method-records for ERM, CDC,
Bloop-style shadow-target rescue, and PCGrad-style shadow-target rescue.  The two
named comparators are non-canonical (\code{canonical:false}); no canonical parity
claim is made.

The frozen-tangent sequence is intentionally fail-closed.  E19 sealed predictions
for five endpoints before training on 16 architecture--seed records.  After that
broad rule failed, E20 froze only drift- and response-crossing labels as primary on
a fresh $4\times8$ crossed data/model-seed factorial per architecture.  E20 remains
the submission-primary empirical study.  Its 10,000-replicate two-way bootstrap
resamples the two observed seed axes.  Surviving manifests, predictions, scores,
hashes, and seals make E19/E20 auditable, but their exact historical dirty source
snapshots were never archived; they are not replayable, and regeneration is a new
study.

\section{Results}
\paragraph{Paired deficits and causal diagnoses.}
At lags 0, 2, and 4, E1 has a positive paired weak-AUC deficit in all 32
$\rho$--seed cells per lag, and all corresponding weak-only controls pass the
specified learnability gate; mean deficits are 5.76, 8.93, and 6.11.  At every
fixed lag, the mean deficit increases over the measured $\rho$ grid.  All 32 lag-8
cells fail the weak-only gate and are weak-only nonlearnable and causally
indeterminate under $\beta=0.25,H=5$; they are neither causal-starvation
certificates nor evidence of absence.  At E-NL's fixed
$(\rho,\mathrm{lag})=(4,2)$ cell, tanh has complete transfer-to-starvation
certificates in 3/8 observed seeds.  GRU has drift and response crossings in 5/8
but 0/8 causal certificates because every weak-only path fails $\beta=0.5$.
Thus rate/outcome suppression differs from causal starvation.  These observed
counts support neither an architecture ranking nor a prevalence claim.

\begin{table}[t]
\centering
\caption{Final nonlinear and frozen-tangent audits. Counts describe the observed
synthetic records, not population prevalence. E20 confidence intervals resample
both seed axes.}
\label{tab:iclr-audit}
\small
\begin{tabular}{llccc}
\toprule
Study & Endpoint & Overall & tanh & GRU \\
\midrule
E-NL & causal certificate & 3/16 & 3/8 & 0/8 \\
E-NL & response crossing & 8/16 & 3/8 & 5/8 \\
E19 & phase / drift / response & 12/16 / 14/16 / 14/16 & --- & --- \\
E19 & causal / learnability & 14/16 / 12/16 & --- & --- \\
E20 & drift event & 60/64, CI $[.8125,1]$ & 32/32 & 28/32 \\
E20 & response event & 64/64, CI $[1,1]$ & 32/32 & 32/32 \\
\bottomrule
\end{tabular}
\end{table}

\paragraph{Gate-indeterminate breadth and comparator-restricted control.}
E26 is negative/indeterminate: MNIST 3-vs-8 and FashionMNIST 0-vs-6 each have
0/32 weak-only-learnable and 0/32 causal records.  The failed gate rejects a false-
positive causal-starvation interpretation, but it also means that zero certificates
cannot establish absence of starvation or intrinsic dataset unlearnability.  E27's
all-$\beta$ architecture endpoint is false for both architectures and was not used
for ranking.  Table~\ref{tab:iclr-breadth} reports the exact macro estimates.  For
E28, weak rescue is $-\lvert\text{weak AUC gap}\rvert$ (higher is better), so
``CDC$-$comparator'' positive favors CDC; trajectory and final deviation are
absolute deviations from paired ERM (smaller is better), so
``comparator$-$CDC'' positive favors CDC.  The frozen joint tradeoff is true only
versus the non-canonical Bloop-style shadow-target rescue.  It is false versus the
non-canonical PCGrad-style shadow-target rescue, so CDC does not beat that variant
under the frozen joint criterion.  No SOTA, canonical parity, or broad CDC
superiority follows.

\begin{table}[t]
\centering
\caption{Clean-source E26--E28 completed outcomes.  Brackets are 95\% CIs; E27
Holm-adjusted $p=1$ for each architecture.}
\label{tab:iclr-breadth}
\small
\resizebox{\linewidth}{!}{%
\begin{tabular}{lll}
\toprule
Study & Endpoint / contrast & Estimate $[95\%\ \mathrm{CI}]$ \\
\midrule
E26 & MNIST weak AUC gap & $-0.0012362842136667493\ [-0.0034405428466004646,0.0010159591371180453]$ \\
E26 & FashionMNIST weak AUC gap & $-0.00010921728159018996\ [-0.004129384520886106,0.003182670049955049]$ \\
E27 & tanh / GRU all-$\beta$ & $0.2916666666666667\ [0.16666666666666666,0.4166666666666667]$ / $0\ [0,0]$ \\
E28 & CDC$-$Bloop-style shadow-target rescue weak rescue & $0.002198259399210656\ [0.0009779187294930126,0.003277366267916477]$ \\
E28 & Bloop-style shadow-target rescue$-$CDC trajectory / final deviation & $3.4651361294978416\ [3.3443827215131754,3.5849715262780952]$ / $0.6422362388111651\ [0.6288088704226539,0.6563168084365315]$ \\
E28 & CDC$-$PCGrad-style shadow-target rescue weak rescue & $-0.00015351238407674823\ [-0.00047909348592838785,0.00013248012419808214]$ \\
E28 & PCGrad-style shadow-target rescue$-$CDC trajectory / final deviation & $-0.48709889128076556\ [-0.6897378944428417,-0.27379327373499074]$ / $-0.2565436437726021\ [-0.3016841153614223,-0.2087651835754514]$ \\
\bottomrule
\end{tabular}}
\end{table}

\paragraph{$\beta$ sensitivity without retraining.}
A post-hoc audit computes each completed weak-only terminal response
$B_{W,s}(H)$ and the fraction above candidate thresholds.  Tanh terminal
minimum, mean, and maximum are $4.41313266754$, $4.49664855003$, and
$4.67605876923$ (8/8 above 0.5); the corresponding GRU values are
$0.0828229486942$, $0.129690139554$, and $0.187142759562$ (0/8).  Thus GRU is not
barely below the selected operating point.  This terminal profile is not generally
equivalent to first-hit learnability on a nonmonotone path and supports neither a
new confirmatory endpoint nor an architecture ranking.

\artifactfigure{../paper/artifacts/enl_tanh_crossover-20260824-141207/enl_learnability_profile.pdf}
{Post-hoc weak-only terminal-response audit from the completed E-NL run. Points
show all eight seeds, horizontal segments show medians, and the dashed line marks
the preregistered $\beta=0.5$. The right panel sweeps candidate thresholds without
rerunning training.}{fig:enl-learnability-profile}

\paragraph{E19 falsification and E20 restricted event-label result.}
E19 obtained phase/drift/response/causal/learnability correctness
$12/16,14/16,14/16,14/16,12/16$, but failed because tanh learnability was
$4/8<5/8$.  Its tanh response-gap, weak-response, and drift RMSE were
$2.367$, $2.394$, and $0.486$; the corresponding GRU values were $0.0163$,
$0.0281$, and $0.0107$.  E20 later passed only the restricted event rule.  Its
secondary phase/causal/learnability correctness was $52/64,56/64,49/64$ (tanh
learnability 17/32); drift/response crossing-time bias was $-0.822/-1.189$; and
tanh response-gap, weak-response, and drift RMSE were $2.455$, $2.489$, and
$0.490$.  Hence E20 is an uncalibrated one-cell event classifier, not
phase, causal, learnability, time, magnitude, prevalence, or kernel-stability
validation.

\paragraph{Historical E3 local mechanism-control ablation.}
In historical E3, all five shadow-based mitigation arms reduce the dense weak AUC
gap from 10.5155 to approximately zero.  CDC alone preserves instantaneous strong
drift to $4.77\times10^{-7}$, but it does not have the best final strong response.
This supports the fixed-chart local mechanism-level guarantee in that historical
synthetic setting, not SOTA performance or broad practical superiority; E28
separately passes the frozen joint tradeoff only against the non-canonical Bloop-
style shadow-target rescue and not against the non-canonical PCGrad-style shadow-
target rescue.

\section{Limitations and conclusion}
Weak learning failure alone is not gradient starvation.  The matched intervention
first asks whether the weak-only path learns under the specified target and
horizon; a failed gate makes causal status indeterminate rather than proving either
starvation or its absence.  Conditional on that premise, rate suppression becomes
outcome suppression only when accumulated negative relative drift overtakes
accumulated positive drift.  The paper's strongest chain is therefore exact finite-
width dynamics $\to$ matched causal intervention $\to$ pathwise attribution
$\to$ local mechanism-level control.  Proposition~\ref{prop:cross-kernel} and the
cumulative balance are deliberately elementary; their value is the exact finite-
width scope and the separation of weak learning failure, rate, outcome, and causal
claims.  Realized $d(0)$ and proved local jets give local or assumption-scoped
signs, while general positive-lag prediction from $\rho$, lag, and recurrent
geometry remains open.

The empirical ceiling is explicit: controlled synthetic task families plus one
semi-real generated-cue intervention, limited fixed cells and seed grids, and no
estimate of universal event prevalence.  E19 remains failed, and E20 remains the
submission-primary empirical study with only its restricted frozen drift/response-
crossing classification endpoints primary.  E26 is not a natural shortcut
benchmark, and 0/32 weak-only learnability per dataset makes its zero causal count
indeterminate rather than evidence of no starvation or intrinsic dataset
unlearnability.  E27's beta-robust endpoint is negative for both architectures and
supports no ranking.  E28's comparators are non-canonical; only the Bloop-style
shadow-target rescue satisfies the frozen joint tradeoff, while CDC does not beat
the PCGrad-style shadow-target rescue under that criterion, so neither SOTA,
canonical parity, nor broad CDC superiority follows.  The unrun Waterbirds adapter
remains a minibatch head-coordinate surrogate outside the full-batch matched-
shadow theorem, and its intercept is not an identified bird-shape response; we make
no Waterbirds or natural-spurious-correlation claim.  E19/E20 are integrity-
auditable but historically unreplayable, whereas E26--E28 used clean
reconstructible source at commit
\code{9ebbc617a9aaad5c3c19e9d5fd4fa08f2d3517d0}.  Their compact archives omit 192
and 1,792 \path{records/**} files, respectively, and therefore are auditable but
not standalone replay packages.
\label{page:main-end}

\subsection*{Reproducibility statement}
The appendix states all theorem assumptions and proofs; gives synthetic and
semi-real data, cue-generation, optimization, seed, censoring, and inference
protocols; maps claims to artifacts and executable guards; and reports failed,
negative, indeterminate, mixed, and successful studies.  Resolved configurations,
environment locks, selection and initialization manifests, authorization seals,
tracked figures, compact outcomes, a dependency-free artifact verifier, and the
full test suite are provided.  The supplement distinguishes reconstructible
clean-source E26--E28, auditable but historically unreplayable E19/E20, and compact
archives that omit record bytes and thus are not standalone replay packages.

\subsection*{AI use statement}
Generative AI tools were used to assist with code drafting and refactoring,
manuscript editing and organization, and review-item tracking.  They were not used
to generate or infer experimental outcomes: every reported number was checked
against saved artifacts or frozen evaluation records.  The authors reviewed the
AI-assisted text and code, ran the test and artifact-validation suites, checked
claim scope against the theorem and evidence ledgers, and take responsibility for
the final content, including all AI-assisted text, claims, code, and artifacts.

\bibliographystyle{plainnat}
\bibliography{references}
\appendix
\fi

\section{Paired causal diagnosis and definitions}
\subsection{Paired logistic gradient flow}
For condition $C\in\{\Bcond,\Wcond\}$, let $\theta_C(\tau)\in\R^P$ follow
full-batch gradient flow on signed logits $r_i^C(\theta)=y_i f_C(x_i;\theta)$:
\begin{equation}
  L_C(\theta)=\frac1{n_C}\sum_{i=1}^{n_C}\log(1+e^{-r_i^C(\theta)}),
  \qquad \dot\theta_C=-\nabla L_C(\theta_C).
  \label{eq:gradient-flow}
\end{equation}
The two conditions share initialization and all data except that the strong channel
is zeroed under $\Wcond$.  This matched weak-only path tests whether the weak signal
can meet the specified target without the strong cue; it is the causal
identification premise that distinguishes suppression from an unsupported weak-
learning label.  Let $M_s^C,M_w^C$ be common differentiable strong and weak
response functionals.  Shared initialization gives
\begin{equation}
  \Delta(0)=0,\qquad
  \Delta(\tau)=M_w^{\Bcond}(\tau)-M_w^{\Wcond}(\tau).
  \label{eq:gap}
\end{equation}
Optimization time $\tau$ is distinct from sequence index $k$.

The controlled dense linear RNN is
\begin{equation}
 h_{k+1}=Wh_k+Bx_k,\qquad f(x)=c^\top h_T,
 \label{eq:linear-rnn}
\end{equation}
with $W_{ij}(0)\sim\mathcal N(0,g^2/N)$,
$B_{ia}(0)\sim\mathcal N(0,1/N)$, and
$c_i(0)\sim\mathcal N(0,1/N)$.  The width scaling of $B$ is essential for a
non-divergent projected metric.  We additionally study tanh RNNs and GRUs.  For
nonlinear models the two responses are common odd symmetric unit-probe contrasts,
$[f(x_a^+)-f(x_a^-)]/2$, not assumed additive coordinates of every data logit.

\subsection{Weak learning failure, causal labels, and metrics}
Let
\begin{equation}
 d(\tau)=\Delta'(\tau)=F_w^{\Bcond}(\tau)-F_w^{\Wcond}(\tau).
 \label{eq:equal-time-drift}
\end{equation}
We call $d>0$ relative transfer and $d<0$ relative rate suppression.  We call
$\Delta<0$ \emph{outcome suppression}, whether it follows a positive excursion or
is present from the start of optimization.  A response crossing is the special case
in which $\Delta$ returns to zero after a positive excursion and then becomes
negative.  Weak-only nonlearnability under the specified gate means that the
weak-only trajectory does not first hit the preregistered target $\beta$ within
horizon $H$.  We reserve \emph{causal starvation} for outcome suppression when that
weak-only trajectory does reach the target at a nondegenerate positive time.  Thus
weak learning failure alone is not gradient starvation.  If the weak-only model
does not reach $\beta$, an outcome crossing is descriptive but not a causal-
starvation certificate; the failed gate also is not evidence that starvation is
absent.  Causal status is indeterminate under the specified target and horizon.

We report first-hitting-time delay when both hits are observed.  Formally,
$\tau_W(\beta;H)=\inf\{\tau\in[0,H]:M_w^{\Wcond}(\tau)\ge\beta\}$, with
$\inf\varnothing=\infty$; the learnability gate requires
$\tau_W(\beta;H)\in(0,H]$, excluding targets met at initialization.  The
horizon-wide weak AUC gap is
\begin{equation}
  \mathcal A_w=\int_0^H
  \left(M_w^{\Wcond}(\tau)-M_w^{\Bcond}(\tau)\right)\dd\tau,
  \label{eq:auc-gap}
\end{equation}
computed by the logged-grid trapezoid rule.  To expose rather than hide sensitivity
to the single preregistered $\beta$, Section~\ref{sec:results} also reports the
threshold-free terminal response $B_W(H)=M_w^{\Wcond}(H)$ across seeds and its
post-hoc threshold profile.  Terminal thresholding is not generally equivalent to
the preregistered first-hit event on a nonmonotone trajectory, so the profile is a
sensitivity diagnostic rather than a replacement gate.  GSI-5 remains a margin-
weight diagnostic, not a definition of starvation.

\section{Exact finite-width dynamics for the paired diagnosis}
The mechanism-level decomposition is exact before any width limit: with
$q_i(r)=\sigm(-r_i)$ and response-specific geometry $c_w=J\nabla M_w$, the weak
response obeys $\dot M_w=c_w^\top q(r)/n$.  The vector $q(r)$ records
cross-entropy weighting across signed sample logits, while $c_w$ records how the
chosen weak response couples to those logits through parameter space.  No DMFT or
two-response closure is needed for this identity; such closures arise only in
stronger predictive approximations.

\begin{proposition}[Universal response--logit cross-kernel flow]
\label{prop:cross-kernel}
At every finite width and every optimization time at which $M_a^C$ is
differentiable,
\begin{equation}
 \dot M_a^C=\frac1{n_C}\sum_i K_{ai}^C\sigm(-r_i^C),
 \qquad K_{ai}^C=\ip{\nabla M_a^C}{\nabla r_i^C}.
 \label{eq:cross-kernel}
\end{equation}
No width limit or two-mode realization is required.  We label this elementary
chain-rule identity as a proposition: its importance here is the exact finite-width
scope and the causal separation it enables, not proof difficulty.
\end{proposition}

\begin{corollary}[Closed projected flow under exact realization]
\label{cor:projected}
If the signed logits admit the exact parameter-independent representation
$r_i^C=\sum_b z_{ib}^C M_b^C$, then
\begin{equation}
 \dot M_a^C=\sum_bG_{ab}^C g_b^C,\quad
 G_{ab}^C=\ip{\nabla M_a^C}{\nabla M_b^C},\quad
 g_b^C=\frac1{n_C}\sum_i z_{ib}^C\sigm(-z_i^C\!\cdot M^C).
 \label{eq:Gg}
\end{equation}
If $r_i^C=z_i^C\cdot M^C+r_i^\perp$, the exact equation instead contains
\begin{equation}
 R_a^C=\frac1{n_C}\sum_i\sigm(-r_i^C)
       \ip{\nabla M_a^C}{\nabla r_i^\perp},
 \qquad \dot M_a^C=(G^Cg^C)_a+R_a^C.
 \label{eq:projection-residual}
\end{equation}
\end{corollary}

\begin{corollary}[Exact full empirical-NTK dynamics]
\label{cor:full-empirical-ntk}
Let $J^C(\tau)\in\R^{n_C\times P}$ have row
$\nabla r_i^C(\theta_C(\tau))^\top$, and define
\begin{equation}
 \Theta^C(\tau)=J^C(\tau)J^C(\tau)^\top,
 \qquad c_a^C(\tau)=J^C(\tau)\nabla M_a^C(\theta_C(\tau)).
\end{equation}
Writing $q(r)=\sigm(-r)$ componentwise, the exact finite-width flow is
\begin{equation}
 \dot r^C=\frac1{n_C}\Theta^C(\tau)q(r^C),
 \qquad
 \dot M_a^C=\frac1{n_C}c_a^C(\tau)^\top q(r^C).
 \label{eq:full-ntk-flow}
\end{equation}
No width limit or response-coordinate realization is required.
\end{corollary}

\begin{definition}[Initialization-frozen nonlinear logistic surrogate]
\label{def:frozen-surrogate}
For one condition, let $\Theta_0=\Theta(0)$ and $c_0=c_a(0)$.  Define
\begin{equation}
 \dot{\bar r}=\frac1n\Theta_0q(\bar r),\qquad
 \dot{\bar M}=\frac1n c_0^\top q(\bar r),\qquad
 (\bar r(0),\bar M(0))=(r(0),M(0)).
 \label{eq:frozen-surrogate}
\end{equation}
The gate remains nonlinear.  This system is exact for the initialization-linearized
tangent model, not for the trained nonlinear network.
\end{definition}

\begin{theorem}[Conditional continuous and secant-SGD comparison]
\label{thm:frozen-comparison}
Use Euclidean and induced operator norms.  Suppose on $[0,t]$ that
$\norm{\Theta(u)-\Theta_0}\le B_\Theta$,
$\norm{c(u)-c_0}\le B_c$, and both logistic weight vectors have norm at most $V$.
Set $a=\norm{\Theta_0}/(4n)$,
\begin{equation}
 \phi_a(t)=\frac{e^{at}-1}{a},\qquad
 \psi_a(t)=\frac{e^{at}-1-at}{a^2},
\end{equation}
with $\phi_0(t)=t$ and $\psi_0(t)=t^2/2$.  Then
\begin{align}
 \norm{r(t)-\bar r(t)}
 &\le \frac{B_\Theta V}{n}\phi_a(t),\label{eq:frozen-logit-bound}\\
 |M(t)-\bar M(t)|
 &\le \frac{B_cV}{n}t
 +\frac{\norm{c_0}B_\Theta V}{4n^2}\psi_a(t).
 \label{eq:frozen-response-bound}
\end{align}

For explicit full-batch SGD, let $\delta\theta_k=\theta_{k+1}-\theta_k$ and define
exact segment kernels
\begin{align}
 \widetilde J_k&=\int_0^1J(\theta_k+s\delta\theta_k)\dd s,&
 \widetilde\Theta_k&=\widetilde J_kJ(\theta_k)^\top,\\
 \widetilde m_k&=\int_0^1\nabla M(\theta_k+s\delta\theta_k)\dd s,&
 \widetilde c_k&=J(\theta_k)\widetilde m_k.
\end{align}
The true increments are exactly
$r_{k+1}-r_k=\eta\widetilde\Theta_kq(r_k)/n$ and
$M_{k+1}-M_k=\eta\widetilde c_k^\top q(r_k)/n$.  If $\eta>0$ and, for every
$0\le j<k$,
$\norm{\widetilde\Theta_j-\Theta_0}\le D_\Theta$,
$\norm{\widetilde c_j-c_0}\le D_c$, and
$\norm{q(r_j)},\norm{q(\bar r_j)}\le V$, then the explicit-Euler
version of Eq.~\eqref{eq:frozen-surrogate} obeys
\begin{align}
 \norm{r_k-\bar r_k}&\le\frac{D_\Theta V}{n}\Phi_k,\\
 |M_k-\bar M_k|&\le\frac{k\eta D_cV}{n}
 +\frac{\eta\norm{c_0}D_\Theta V}{4n^2}\Psi_k,
 \label{eq:frozen-discrete-bound}
\end{align}
where
\begin{equation}
 \Phi_k=\frac{(1+\eta a)^k-1}{a},\qquad
 \Psi_k=\frac{(1+\eta a)^k-1}{\eta a^2}-\frac{k}{a},
\end{equation}
with limits $\Phi_k=k\eta$ and $\Psi_k=\eta k(k-1)/2$ at $a=0$.
Apply the bounds separately to $\Bcond,\Wcond$ and add them to bound the paired
response gap.  Frozen signs or brackets transfer only when their strict margins
exceed the corresponding summed errors; first hitting times additionally require
Corollary~\ref{cor:hitting}'s prehistory, transversality, radius, and mesh
hypotheses.
\end{theorem}

Theorem~\ref{thm:frozen-comparison} is conditional: this paper does not prove the
required instantaneous or secant-kernel movement bounds for tanh or GRU.  A held-
out event-classification result cannot substitute for that premise.  In particular,
the theorem is finite-width and is not a DMFT closure.

For channel-localized modes of the dense linear RNN, differentiation across the
disjoint $(c,B,W)$ blocks gives the exact finite-width geometry
\begin{equation}
 G_{ab}=p_a^\top p_b+
 \ind\{\operatorname{ch}(a)=\operatorname{ch}(b)\}q_a^\top q_b+
 \ip{S_a}{S_b}_F,
 \label{eq:dense-geometry}
\end{equation}
where $p_a=W^{\ell_a}b_a$, $q_a=(W^{\ell_a})^\top c$, and
\begin{equation}
 S_a=\sum_{j=0}^{\ell_a-1}(W^j)^\top c
                 (W^{\ell_a-1-j}b_a)^\top.
\end{equation}
\begin{theorem}[Exact realized-initialization rate-gap criterion]
\label{thm:initial-gap}
In the noiseless positive rank-one dense-linear pair, let the strong cue have
strength $\rho>0$ at lag zero, the unit weak cue have lag $\ell$, and define
\begin{align}
 K&=G_{ws}(0)=b_s^\top W^\ell b_w,\\
 A&=G_{ww}(0)=\norm{W^\ell b_w}^2+\norm{(W^\ell)^\top c}^2+\norm{S_w}_F^2,\\
 S_w&=\sum_{j=0}^{\ell-1}(W^j)^\top c
       (W^{\ell-1-j}b_w)^\top,\\
 s_B&=\sigm[-(\rho M_s(0)+M_w(0))],\qquad
 s_W=\sigm[-M_w(0)].
\end{align}
(The recurrent sum is zero for $\ell=0$.)  The exact initial equal-time weak-rate
gap is
\begin{equation}
 d_0=s_B(\rho K+A)-s_WA.
 \label{eq:initial-gap-certificate}
\end{equation}
Consequently, initial transfer, suppression, and the boundary occur respectively
when
\begin{equation}
 K\gtrless\frac{A}{\rho}\left(\frac{s_W}{s_B}-1\right)
 \quad\text{or equality holds.}
 \label{eq:initial-gap-threshold}
\end{equation}
This is an exact pre-trajectory criterion from the realized initialization.  Since
shared initialization gives $\Delta(0)=0$ and $\Delta'(0)=d_0$, it also implies
\begin{equation}
 d_0<0\Longrightarrow \Delta(t)<0,
 \qquad
 d_0>0\Longrightarrow \Delta(t)>0
 \quad\text{for all sufficiently small }t>0.
 \label{eq:initial-gap-local-outcome}
\end{equation}
These strict local outcome signs are existential and provide no quantitative
horizon.  A numerical $d_0$ within the executable certificate's caller tolerance
is reported as boundary/undetermined and licenses neither sign.  The criterion is
not a distribution-only positive-lag phase boundary: under the stated iid scaling,
$K$ and $M_s(0)$ contribute a seed-dependent leading sign.
\end{theorem}

Theorem~\ref{thm:initial-gap} is a partial predictor: realized $d(0)$ gives an
exact local sign in its stated dense-linear scope, not a global trajectory or a
parameter-only law in $\rho$ and lag.  Where a higher local jet is proved---as in
the zero-lag result with $d(0)=0$ and $d'(0)<0$---it likewise supplies a rigorous
local sign only in that stated scope.  Stronger positive-lag trajectory prediction
requires independently established product-ball or kernel-movement control, for
which this paper supplies no nonvacuous general constants.

The proofs are in Appendix~\ref{app:finite-proofs}.  Equation~\eqref{eq:Gg} is
exact for the zero-background-noise linear synthetic model.  For tanh/GRU probes,
we use Proposition~\ref{prop:cross-kernel}'s direct autograd drift as primary and
report $R$; projected and matched-state decompositions are diagnostics only.

\section{Separating rate suppression, outcome suppression, and causal starvation}
\begin{proposition}[Cumulative drift balance]
\label{prop:drift-balance}
Let $\Delta$ be absolutely continuous on $[0,H]$, with $\Delta(0)=0$ and
$d=\Delta'$ almost everywhere.  Write
\begin{equation}
 d_+(t)=[d(t)]_+,
 \qquad d_-(t)=[-d(t)]_+,
 \qquad
 P(t)=\int_0^t d_+(u)\dd u,
 \qquad N(t)=\int_0^t d_-(u)\dd u.
\end{equation}
Then $\Delta(t)=P(t)-N(t)$ for every $t\in[0,H]$.  Thus strict outcome
suppression occurs by $H$ if and only if $N(t)>P(t)$ for some $t\in(0,H]$;
this statement permits arbitrary sign changes and suppression from initialization.
This is deliberately a bookkeeping proposition: its contribution is the causal
separation it enforces, not proof difficulty.
\end{proposition}

\begin{corollary}[Single-crossing tail-area criterion]
\label{cor:tail-area}
Suppose additionally that there is $\tau_d\in(0,H)$ such that $d>0$ almost
everywhere on $(0,\tau_d)$ and $d<0$ almost everywhere on $(\tau_d,H)$.  Define
\begin{equation}
 A_+=P(\tau_d)=\Delta(\tau_d)>0.
\end{equation}
Then $P(t)=A_+$ and $\Delta(t)=A_+-N(t)$ for $t\ge\tau_d$.  Hence a later
equality is reached by $H$ if and only if $N(H)\ge A_+$, and strict outcome
suppression occurs by $H$ if and only if $N(H)>A_+$.  If $d<0$ on each
nontrivial post-crossing interval, the later equality is unique whenever it exists.
\end{corollary}

A $+\to-$ drift crossing therefore locates a local maximum of $\Delta$; it does
not force $\Delta$ to return to zero.  Pathwise attribution requires the negative
area to repay all positive area already accumulated: outcome suppression begins
only when accumulated suppression overtakes accumulated help.  Causal
\emph{starvation} further requires the independent weak-only learnability gate.  A
failed gate blocks that certificate but does not establish that starvation is
absent.

\begin{corollary}[Response-crossing bounds]
\label{cor:crossing-bounds}
Under Corollary~\ref{cor:tail-area}, if $d(t)\le-\kappa<0$ on $[t_0,H]$ with
$t_0>\tau_d$, then, when the right side lies within the horizon,
\begin{equation}
 \tau_{\mathrm{response}}\le
 t_0+\frac{[A_+-N(t_0)]_+}{\kappa}
 \le t_0+\frac{A_+}{\kappa}.
\end{equation}
If $d(\tau_d)=0$ and $d'(t)\le-\lambda<0$ on
$[\tau_d,\tau_d+r]$, then
$N(t)\ge\lambda(t-\tau_d)^2/2$; when
$A_+\le\lambda r^2/2$,
\begin{equation}
 \tau_{\mathrm{response}}\le\tau_d+\sqrt{2A_+/\lambda}.
\end{equation}
\end{corollary}

\begin{theorem}[Conditional parameter-tube certificate]
\label{thm:parameter-tube}
Let $X=(\theta_B,\theta_W)$, let $V$ be the joint paired gradient-flow field,
and define $d=L_V\Delta$ and $j=L_Vd$.  On the radius-$R$ product-norm ball
$U_R$ about the shared initialization, suppose uniformly that
\begin{equation}
 \norm{V(X)}\le S_R,\qquad
 -\Lambda\le j(X)\le-\lambda<0,
 \quad 0<\lambda\le\Lambda,
 \label{eq:parameter-tube-premises}
\end{equation}
with $S_R>0$, and set $h_R=R/S_R$.  If $d_0>0$ and
$h_R>d_0/\lambda$, there is exactly one transfer-to-rate-suppression crossing
before ball exit and
\begin{equation}
 \frac{d_0}{\Lambda}\le\tau_d\le\frac{d_0}{\lambda}.
 \label{eq:parameter-tube-crossing}
\end{equation}
If $h_R>2d_0/\lambda$, then $\Delta(t)<0$ for every
$2d_0/\lambda<t\le h_R$.  If also the weak-only drift obeys
$\dot M_w^{\Wcond}\ge\nu>0$ uniformly, the target is nondegenerate, and
\begin{equation}
 \beta>M_w^{\Wcond}(0),\qquad
 \frac{\beta-M_w^{\Wcond}(0)}{\nu}<h_R,
 \label{eq:parameter-tube-learnability}
\end{equation}
then the weak-only member learns by reaching $\beta$ at a strictly positive time
within the same certified horizon, so the outcome suppression is causally gated
starvation rather than suppression alone.  If
$\beta\le M_w^{\Wcond}(0)$, the target was met at initialization; this is
degenerate attainment, not learnability, and cannot certify causal starvation.
\end{theorem}

\begin{corollary}[Certified finite-horizon safe transfer]
\label{cor:parameter-safe}
On the same ball, if instead $|j(X)|\le J$ and $d_0-Jh_R>0$, then for
$0<t\le h_R$,
\begin{equation}
 d(t)\ge d_0-Jt>0,
 \qquad
 \Delta(t)\ge d_0t-\frac{Jt^2}{2}>0.
 \label{eq:parameter-safe}
\end{equation}
Thus suppression is excluded on this certified finite horizon, not for all future
optimization time.  The executable checker fails closed when optional premises are
missing; it does not estimate $S_R,\lambda,\Lambda,\nu$, or $J$ from sampled
trajectories.  No current component constructs nonvacuous general positive-lag
constants.
\end{corollary}

\subsection{An ordering-invariant rank-one criterion}
Assume deterministic positive cues $z^{\Bcond}=(\rho,1)$,
$z^{\Wcond}=(0,1)$, $\rho>0$, and exact two-mode realization.  Define
\begin{align}
 u_B&=\rho M_s^{\Bcond}+M_w^{\Bcond},& s_B&=\sigm(-u_B),\\
 u_W&=M_w^{\Wcond},&s_W&=\sigm(-u_W),\\
 H_B&=\rho G_{ws}^{\Bcond}+G_{ww}^{\Bcond},&
 H_W&=G_{ww}^{\Wcond},\\
 Q_B&=\rho^2G_{ss}^{\Bcond}+2\rho G_{sw}^{\Bcond}+G_{ww}^{\Bcond}\ge0.
\end{align}

\begin{theorem}[Log drift-ratio identity]
\label{thm:log-ratio}
The weak drifts factor as $F_w^{\Bcond}=s_BH_B$ and
$F_w^{\Wcond}=s_WH_W$.  Where $H_B,H_W>0$, define
\begin{equation}
 \Psi=\log\frac{F_w^{\Bcond}}{F_w^{\Wcond}}
 =\log\frac{s_B}{s_W}+\log\frac{H_B}{H_W}.
\end{equation}
Then $\operatorname{sign}(d)=\operatorname{sign}(\Psi)$ and
\begin{equation}
 \Psi'=-s_B(1-s_B)Q_B+s_W(1-s_W)H_W
       +\left(\log\frac{H_B}{H_W}\right)'.
 \label{eq:Psi-prime}
\end{equation}
\end{theorem}

\begin{corollary}[Unique rate crossing]
\label{cor:unique-rate-crossing}
Suppose on $[t_0,H]$ that $H_B,H_W>0$, $\Psi(t_0)>0$,
$\Psi'(t)\le-\kappa<0$, and $H-t_0\ge\Psi(t_0)/\kappa$.  Then exactly one
rate crossing occurs in $(t_0,H]$ and
\begin{equation}
 \tau_d\le t_0+\Psi(t_0)/\kappa.
\end{equation}
A sufficient primitive condition is
\begin{equation}
 s_B(1-s_B)Q_B\ge a,\quad s_W(1-s_W)H_W\le b,\quad
 (\log(H_B/H_W))'\le c,\quad a-b-c\ge\kappa>0.
\end{equation}
\end{corollary}

The bounds must be established independently to be predictive.  Evaluating
$\Psi$ on a completed trajectory is certification, not prediction.  Noisy cues
generally break the rank-one factorization, and even a rate crossing still requires
Corollary~\ref{cor:tail-area} for an outcome crossing.

\subsection{Hitting-time stability}
\begin{corollary}[Transverse first-hit stability]
\label{cor:hitting}
Let continuous $x_N,x:[0,H]\to\R$ satisfy
$\norm{x_N-x}_\infty\le\epsilon$, and define
\begin{equation}
 \tau=\inf\{t\in[0,H]:x(t)\ge\beta\},\qquad
 \tau_N=\inf\{t\in[0,H]:x_N(t)\ge\beta\},
 \quad \inf\varnothing=\infty.
\end{equation}
Assume $x$ is absolutely continuous, $x(0)<\beta$, and
$\tau\in(0,H)$.  Choose $0<r\le\min\{\tau,H-\tau\}$ such that
$x'(t)\ge\kappa>0$ almost everywhere on $[\tau-r,\tau+r]$, and let
\begin{equation}
 \eta=\beta-\max_{0\le t\le\tau-r}x(t)>0.
\end{equation}
If $\epsilon<\eta$ and $\epsilon/\kappa\le r$, then $\tau_N<\infty$ and
\begin{equation}
 |\tau_N-\tau|\le\epsilon/\kappa.
 \label{eq:hitting-continuous}
\end{equation}
For an ordered grid $0=t_0<\cdots<t_m=H$ of maximum mesh
$h=\max_j(t_j-t_{j-1})$, let $\tau_N^{\rm samp}$ be the first sampled time with
$x_N(t_j)\ge\beta$, and let $\tau_N^{\rm int}$ be the first hit of the
piecewise-linear interpolant through the samples.  If
$\epsilon/\kappa+h\le r$, both are finite and
\begin{equation}
 |\tau_N^{\rm samp}-\tau|\le\epsilon/\kappa+h,
 \qquad
 |\tau_N^{\rm int}-\tau|\le\epsilon/\kappa+h.
 \label{eq:hitting-grid}
\end{equation}
For a paired delay, if both reference hits are finite and each trajectory satisfies
its own assumptions, then
\begin{equation}
 |(\tau_{N,B}-\tau_{N,W})-(\tau_B-\tau_W)|
 \le \epsilon_B/\kappa_B+\epsilon_W/\kappa_W+h_B+h_W.
\end{equation}
If either hit is infinite, the delay is censored; a common mesh $h$ gives the
special case with $2h$.
\end{corollary}
Descending crossings follow by sign reversal and localization after a strict
excursion.  Since the response gap has $\Delta(0)=0$, its later downward outcome
crossover must be localized to $[a,H]$ with $\Delta(a)>0$ and treated using
$x=-\Delta$ and target zero; the hit at time zero is not that later crossover.
This corollary transfers an independently proved uniform trajectory limit.  It
does not prove that limit, first-entry semantics, the prehistory margin, or the
derivative bound.  Thus Theorem~\ref{thm:frozen-comparison}'s summed paired
response bound licenses a crossing-time statement only after all of these separate
guards hold; an error formula or empirical event match alone is not a hitting-time
certificate.

\section{A solvable recurrence-invisible zero-lag wide limit}
Set lag separation zero and background input/noise zero in the dense linear model;
let the realized recurrent matrix $W$ be arbitrary and finite.  Both cues then
arrive at the final step.  Starting from $h_0=0$, induction gives $h_t=0$ for every
state entering the final update and hence
$h_T=Wh_{T-1}+Bx_{T-1}=Bx_{T-1}$.  Thus $W$ is loss-invisible,
$\nabla_WL=0$, and it remains fixed.  Writing the input columns as $b_s,b_w$ gives
$M_a=c\cdot b_a$.  Let $u_{ab}=b_a\cdot b_b$ and $w=c\cdot c$.

\begin{theorem}[Every-width six-scalar closure]
\label{thm:six-scalar}
At every width,
\begin{align}
 \dot M_a&=g_aw+\sum_b g_bu_{ab},&
 \dot u_{ab}&=g_aM_b+g_bM_a,\\
 \dot w&=2\sum_bg_bM_b,&
 G_{ab}&=u_{ab}+\delta_{ab}w.
 \label{eq:six-scalar}
\end{align}
\end{theorem}

\begin{theorem}[Finite-horizon deterministic wide limit]
\label{thm:wide-limit}
For any finite realized recurrent matrix $W$, let all coordinates of
$b_s,b_w,c$ be mutually independent $\mathcal N(0,1/N)$ and
\begin{equation}
 q_N=(M_s,M_w,u_{ss},u_{sw},u_{ww},w),\qquad
 q_0=(0,0,1,0,1,1).
\end{equation}
For every $\epsilon>0$,
\begin{equation}
 \Prob\!\left(\norm{q_N(0)-q_0}_\infty>\epsilon\right)
 \le \min\left\{1,\frac{9}{N\epsilon^2}\right\}.
 \label{eq:init-concentration}
\end{equation}
Fix a cue law with bounded second moments and define
\begin{equation}
 g_a(M)=\E[z_a\sigm(-z\cdot M)],\quad
 A=\max_a\E|z_a|,\quad
 C=\frac14\max_a\sum_b\E|z_az_b|.
\end{equation}
Then $\norm{g(M)}_\infty\le A$ and
\begin{equation}
 \norm{g(M)-g(\widetilde M)}_\infty
 \le C\norm{M-\widetilde M}_\infty.
 \label{eq:g-lipschitz}
\end{equation}
Let $q(t)$ be the deterministic solution from $q_0$, let
$c_\star=W(e^{-1})=\sup_x x\sigm(-x)$, where $W$ is Lambert's function, and set
\begin{equation}
 R_H=3+4c_\star H,\qquad B_{H,r}=R_H+r,\qquad
 L_{H,r}=4(A+B_{H,r}C)
 \label{eq:deterministic-tube-constants}
\end{equation}
for a fixed deterministic radius $r>0$.  The reference Gram energy
$R=w+u_{ss}+u_{ww}$ obeys $R(0)=3$, $R'\le4c_\star$, and therefore
$\norm{q(t)}_\infty\le R_H$ on $[0,H]$.  The vector field is
$L_{H,r}$-Lipschitz in infinity norm on the radius-$r$ tube about this path.
If
\begin{equation}
 \norm{q_N(0)-q_0}_\infty
 \le e^{-L_{H,r}H}\min\{r,\delta\},
 \label{eq:tube-bootstrap-event}
\end{equation}
then $q_N$ does not leave the tube before $H$ and
\begin{equation}
 \sup_{t\le H}\norm{q_N(t)-q(t)}_\infty\le\delta.
 \label{eq:gronwall}
\end{equation}
Consequently,
\begin{equation}
 \Prob\!\left(\sup_{t\le H}\norm{q_N(t)-q(t)}_\infty>\delta\right)
 \le\min\left\{1,
 \frac{9e^{2L_{H,r}H}}{N\min\{r,\delta\}^2}\right\}.
 \label{eq:trajectory-concentration}
\end{equation}
For $r=1$ and $0<\delta\le1$, the denominator reduces to $N\delta^2$ with
the deterministic constant $L_{H,1}$.  Thus $q_N\to q$ uniformly on fixed
horizons in probability, with initialization scale $O_{\Prob}(N^{-1/2})$.
For the paired both/weak-only result, use a common $r$ and the maximum of the two
deterministic $L_{H,r}$ constants for the two fixed cue laws.
\end{theorem}

This theorem is a genuine large-width result at recurrence-invisible zero lag for
any finite realized recurrent gain.  It is not a positive-lag recurrent DMFT.

\begin{corollary}[Zero-lag suppression and weak-only learnability]
\label{cor:zero-disorder-suppression}
For the noiseless positive rank-one cue law
($\texttt{cue\_noise}=0$), zero lag, zero background noise, and any finite
realized recurrent matrix, use the deterministic symmetric initial state
$(M_s,M_w,u_{ss},u_{sw},u_{ww},w)=(0,0,1,0,1,1)$.  For every $\rho>0$, the
both-minus-weak gap and rate satisfy
\begin{equation}
 \bar\Delta(0)=0,\qquad \bar d(0)=0,\qquad
 \bar d'(0)=-\rho^2/2<0.
 \label{eq:zero-disorder-initial-suppression}
\end{equation}
Hence $\bar\Delta(t)<0$ for all sufficiently small $t>0$: suppression begins at
initialization and is not transfer-then-suppression.  In the weak-only subsystem,
\begin{equation}
 u_{ww}=w=\sqrt{1+M_w^2},\qquad
 \dot M_w=2\sigm(-M_w)\sqrt{1+M_w^2}>0,
\end{equation}
so every finite $\beta>0$ is reached and
\begin{equation}
 \tau_\beta\le\frac{\beta}{2\sigm(-\beta)}.
 \label{eq:zero-disorder-target-time}
\end{equation}
These identities and the target bound are independent of the arbitrary finite
realized $W$.  The finite-width outcome-suppression/learnability conclusion
follows from
Theorem~\ref{thm:wide-limit} only after deterministic witness margins and its
common tube event are established; ordinary floating-point integration is not
silently promoted to such a proof.
\end{corollary}

\section{Local mechanism-level control with Counterfactual Drift Correction}
CDC uses the paired weak-only drift target to intervene on the diagnosed local
mechanism; it is not presented as a general optimizer or robustness method.
Flatten the implemented trainable tensors into one fixed parameter chart with the
standard product Euclidean/Frobenius inner product.  Every gradient, dot product,
norm, projection, Lipschitz ball, and straight segment in this section uses that
chart and metric; no reparameterization-invariance or natural-gradient claim is
made.

At a both-feature state $\theta$, let
\begin{equation}
 a=\nabla m_s(\theta),\quad p=\nabla m_w(\theta),\quad
 v=-\nabla L(\theta),\quad F_s=a\cdot v,\quad F_w=p\cdot v.
\end{equation}
A same-initialization weak-only shadow supplies $F_w^{\Wcond}$.  Define
\begin{equation}
 q=\begin{cases}
 p-\dfrac{p\cdot a}{\norm a^2}a,&a\ne0,\\
 p,&a=0,
 \end{cases}
 \quad
 \delta=[F_w^{\Wcond}-F_w]_+,
 \quad
 \alpha=\begin{cases}\delta/\norm q^2,&\delta>0,\ q\ne0,\\0,&\delta=0.
 \end{cases}
 \label{eq:cdc}
\end{equation}
The corrected velocity is $v'=v+\alpha q$.  A positive deficit with $q=0$ is
infeasible.  The theorems assume full-batch gradients, plain SGD, zero weight
decay, and no clipping.  A nonbinding coefficient cap leaves the uncapped result
unchanged; a binding cap preserves orthogonality but generally loses target
attainment, so no target-attaining optimality claim then applies.  In the fixed
implemented chart, feasible uncapped target-attaining CDC is the unique minimum-
Euclidean-norm correction satisfying zero instantaneous strong-drift change and
the weak-drift lower bound.  Its finite-step guarantee is local and one-step.

\begin{theorem}[CDC local guarantees]
\label{thm:cdc}
The construction in Eq.~\eqref{eq:cdc} has the following fixed-coordinate
properties.
\begin{enumerate}[leftmargin=*]
  \item \textbf{Instantaneous preservation:} $a\cdot v'=a\cdot v$ exactly,
  including when $\alpha$ is capped.
  \item \textbf{Minimum Euclidean norm:} for the feasible uncapped problem,
  $d^*=\alpha q$ is the unique minimum-norm, lower-bound-attaining correction
  satisfying $a\cdot d=0$ and
  $p\cdot(v+d)\ge F_w^{\Wcond}$.
  \item \textbf{Local finite-step bound:} if $\nabla m_s$ is $L$-Lipschitz in
  the same norm on a convex neighborhood containing both straight coordinate step
  segments, then
  \begin{equation}
  |m_s(\theta+\eta v')-m_s(\theta+\eta v)|
  \le\frac L2\eta^2(\norm{v'}^2+\norm v^2).
  \label{eq:cdc-step}
  \end{equation}
  For the active uncapped correction, if $\norm v\le V$, $\delta\le D$, and
  $\norm q\ge q_{\min}>0$, the right side is at most
  \begin{equation}
  \frac L2\eta^2\left[V^2+(V+D/q_{\min})^2\right].
  \end{equation}
  This explicit norm bound also remains valid under a nonnegative upper cap,
  which can only reduce $\alpha$.
\end{enumerate}
\end{theorem}
For a positive raw deficit, exact feasibility can be stated intrinsically only as
$dm_w$ having a nonzero restriction to $\ker(dm_s)$; in this metric it is
equivalent to $q\ne0$.  The represented $q$, its norm, the numerical
\code{feasibility\_epsilon}, and the selected minimum-norm correction are
metric- and parameter-scale-dependent.  The finite-step result is one-step and
local in the same chart.  It is not a trajectory guarantee and does not imply
final strong-response retention.

\section{Executable scope guards}
Appendix Table~\ref{tab:contracts} maps theorem objects to their implementations
and explicit scope guards.  Helpers enforce only represented scalar or
configuration conditions; callers still establish trajectory, model, first-entry,
and smoothness hypotheses.  Unit tests check software identities and edge cases,
not mathematical propositions or unrepresented assumptions.

Every general positive-lag solver path that would require an unproved closure
raises \code{NotImplementedError}; arbitrary finite recurrent gain is accepted only
on the proved zero-lag recurrence-invisible path.  This is a scientific guard, not
missing numerical polish.

\section{Experimental protocol}
\subsection{Synthetic task and pairing}
Sequences have two input channels.  A strong cue of scale $\rho$ and a unit weak
cue appear at controlled sequence positions separated by a lag.  Positive tasks
align the cues with the label; the negative control reverses the strong cue.  Data
and initialization are paired across both-feature and weak-only conditions.
Training is deterministic full-batch cross-entropy with zero weight decay.  Except
for E-NL, the two conditions can be run sequentially because deterministic
full-batch updates preserve the same trajectories; E-NL uses lockstep models so
both exact drifts are available at each logged time.

Appendix Table~\ref{tab:protocol} records model, task, seed, and compute
settings, including the roughly doubled per-step model work required by matched-
shadow methods.  E1 uses four strengths and seed IDs 4--11, so each lag contains 32
$\rho$--seed cells but only eight independent seed IDs.  Its gain is $0.5$,
$\beta=0.25$, and horizon is $H=5$; the negative control is one eight-seed cell at
$(\rho,\mathrm{lag})=(4,4)$.  E2 uses gain $0.5$ and a held-out empirical closure
reference.  E-NL uses one cell ($\rho=4$, lag 2), eight seeds per architecture,
gain $0.8$, cue noise $0.1$, zero background noise, $\beta=0.5$, and direct
autograd drift.  The frozen-kernel studies use the same task cell.  Historical E3
compares ERM, Spectral Decoupling, an interaction penalty, CDC, loss-gradient projection,
unconstrained rescue, Bloop-style shadow-target rescue, and PCGrad-style
shadow-target rescue; the five rescue arms share the same weak-only shadow target.
These style labels do not assert canonical source parity.

\paragraph{Prospectively frozen gate and mechanism-control extensions.}
E26 executes 64 crossed records (four data seeds by eight model seeds per dataset)
on MNIST 3-vs-8 and FashionMNIST 0-vs-6.  It retains real core pixels and ablates
only an exactly removable generated cue channel, so its publication scope is
semi-real rather than natural-spurious-correlation evidence.  E27 executes 192
nonlinear records over tanh/GRU, three fixed $(\rho,\mathrm{lag})$ cells, and the
prospectively frozen first-hit grid $\beta\in\{0.25,0.5,1,2\}$.  E28 executes 256
method-records over ERM, CDC, Bloop-style shadow-target rescue, and PCGrad-style
shadow-target rescue.  The latter two are response-gradient shadow-target variants
with \code{canonical:false}; no canonical algorithm/reference parity is claimed.
All three studies were externally sealed and executed from clean commit
\code{9ebbc617a9aaad5c3c19e9d5fd4fa08f2d3517d0}.  The excluded evaluator-
engineering set and its non-scientific use are documented in
Appendix~\ref{app:experiment-audit}.

\paragraph{Manifest-sealed surrogate audit protocol.}
\code{enl-preflight} constructs no optimizer and observes no trained state.  It
writes the initialization kernels and predictions, verifies state/gradient
nonmutation, and hashes every declared file and per-record kernel together with the
resolved config and executable-source fingerprint.  Before constructing an
optimizer, \code{enl-evaluate} requires the sidecar, actual manifest, and externally
pinned digest to agree; verifies every hash/size, source, protocol, seed pair, and
time grid; and loads rather than rebuilds predictions.  It re-hashes all sealed
inputs after training.  The broad five-endpoint rule was frozen before its pilot;
after that rule failed, only drift/response crossing was frozen as primary on fresh
seeds.  The latter crossed design uses a two-way pigeonhole bootstrap with 10,000
replicates.

\paragraph{Statistics and censoring.}
Independent seed pairs are the unit of replication only when neither seed axis is
reused.  E1's 32 cells per lag reuse eight seed IDs across four strengths, and the
E-NL 3/8, 5/8, and 0/8 proportions describe only the eight observed seeds per
architecture; none is an estimate of universal event prevalence.  Under E20's
crossed data/model-seed reuse, row-wise Student-$t$ intervals are suppressed and
the two-way percentile interval resamples the observed seed axes independently.
It is conditional on those observed levels: an all-correct table necessarily gives
$[1,1]$ and does not quantify unseen failure modes.  Hitting times beyond the
horizon are right-censored and never converted into finite delays.  Conditional
crossing-time means are reported only with their crossing-record denominator.  The
exploratory E-NL $5\times5$ grid uses seeds 0--3 and is explicitly post-hoc; it is
not used to select a confirmatory condition.  E26 reports dataset-specific two-way
crossed-seed intervals; absent weak-only first hits censor causal interpretation in
both directions.
E27 uses a shared 5,000-replicate two-way pigeonhole bootstrap, equal-weights the
three fixed cells, and Holm-adjusts the two architecture tests; architecture
ranking was forbidden.  E28 uses the same bootstrap draws and equal-weights its two
fixed cells.  Its weak-rescue contrast is CDC minus comparator for
$-\lvert\text{weak AUC gap}\rvert$ (positive favors CDC); its trajectory and final
contrasts are comparator minus CDC absolute deviation from paired ERM (positive
favors CDC).

\paragraph{Reproducibility.}
Ordinary runs write resolved configuration, environment, trajectory, summary, and
figure artifacts; frozen studies additionally seal predictions and hashes before
training.  The final E-NL run archives a 17-file executable-source fingerprint and
unmodified metadata, making that named run reconstructible.  E19 and E20 are
different: their manifests, predictions, scores, seals, and historical source
fingerprints survive, but the exact dirty Python source bytes were not archived and
are unavailable and unreconstructible.  Current source intentionally fails their
source guard; a new preflight/evaluation is a new study, not a replay.  Historical
E1, E2, E2-R, exploratory E-NL, and E3 artifacts likewise disclose dirty states
predating source fingerprinting.  By contrast, E26--E28 were executed from clean
commit \code{9ebbc617a9aaad5c3c19e9d5fd4fa08f2d3517d0}; their source manifests
make the source reconstructible.  The compact E26 archive omits 192
\path{records/**} files and the shared E27/E28 archive omits 1,792, so hashes and
central outputs remain auditable but neither compact archive is a standalone replay
package.

\section{Results}
\label{sec:results}
\subsection{Paired weak-AUC deficits and weak-only learnability}
At lags 0, 2, and 4, all 32 $\rho$--seed cells per lag (four strengths by
eight seed IDs) have positive paired weak-AUC deficits and learnable weak-only
controls under the specified gate.  Averaged across $\rho$, the lag-wise deficits
are 5.76, 8.93, and 6.11, respectively; the measured effect peaks at lag 2 rather
than decreasing monotonically with lag.  At every fixed lag, the mean deficit
increases monotonically over the measured $\rho$ grid.  The negative-control cell
is transfer in 8/8 seeds, with mean gap $-4.038$ and 95\% CI
$[-4.498,-3.578]$.  Every lag-8 cell fails the weak-only gate: all 32 are weak-only
nonlearnable and causally indeterminate under $\beta=0.25,H=5$.  They supply
neither causal-starvation certificates nor evidence that starvation is absent.
Final accuracy is 1.0 throughout, and no row meets the target at initialization.

\artifactfigure{../paper/artifacts/historical/e1_dense_rerun-20260823-085142/e1_causal_regions.pdf}
{Controlled dense-linear paired weak-AUC deficits and weak-only gate status.  The
specified learnability gate makes all lag-8 points causally indeterminate rather
than positive or negative starvation evidence.}{fig:e1}

\subsection{Empirical width trends and exact zero-lag checks}
E2 compares finite networks with a disjoint-seed empirical
\emph{closure reference}; it is not an independently derived theory.  Overall
both-condition trajectory NRMSE is $0.09084,0.04761,0.05149,0.02958$ at widths
$32,64,128,256$.  Weak-only NRMSE is
$1.56931,0.38792,0.40201,0.20024$.  Both improve end-to-end but are non-monotone at
$N=128$.  All 32 point--condition--diagnostic combinations improve end-to-end;
large normalized weak-only values at difficult points have near-zero denominators.
The exact finite-width geometry check has zero recorded relative error and the
maximum projected-identity error is $2.86\times10^{-6}$.

In the exact legacy-named zero-disorder solver check, the time-discretization
residual scales as
$O(\eta)$ with fitted slope 1.0115.  A frozen-geometry integrator agrees with the
projected-flow integrator to $1.91\times10^{-15}$.  The general solver acceptance
record remains false by construction because all general-DMFT checks are blocked.

\artifactfigure{../paper/artifacts/historical/e2_corrected_validation-20260821-224630/e2_field_agreement.pdf}
{Finite-width agreement with a held-out empirical finite-network closure reference.
This is not DMFT validation.}{fig:e2}

\subsection{Tanh certificates versus ungated GRU crossings}
Table~\ref{tab:enl} is the authoritative nonlinear result at the fixed
$(\rho,\mathrm{lag})=(4,2)$ cell.  Tanh has a complete exact-drift/response/tail-
area transfer-to-starvation certificate in 3/8 observed seeds; the other five have
negative exact relative drift from initialization and therefore do not exhibit the
transfer-to-starvation path type.  GRU has exact drift and response crossings in
5/8 seeds, but every weak-only trajectory fails to reach $\beta=0.5$; its causal
count is therefore 0/8 and causal status is indeterminate under that gate.  This
separates rate/outcome suppression from causal starvation; it does not rank
architectures or estimate prevalence.  The nonzero direct-versus-projected
residuals empirically confirm that $Gg$ cannot be silently treated as the exact
nonlinear probe derivative.

\begin{table}[t]
\centering
\caption{Final E-NL result at $\rho=4$, lag 2, eight seeds.  Crossing times and
confidence intervals are conditional on crossing and must not be read as
eight-seed estimates.}
\label{tab:enl}
\small
\begin{tabular}{lcccccc}
\toprule
Model & exact drift & response & causal & mean $\tau_d$ & mean $\tau_r$ & max $|R_w^B-R_w^W|$ \\
\midrule
tanh & 3/8 & 3/8 & 3/8 & 0.194 & 0.389 & 0.02197 \\
GRU  & 5/8 & 5/8 & 0/8 & 3.846 & 7.200 & 0.002986 \\
\bottomrule
\end{tabular}
\end{table}

The preregistered gate is the first hit at $\beta=0.5$, but a single binary
operating point can hide whether failures are marginal.  We therefore add a post-
hoc terminal audit using $B_{W,s}(H)=M_{w,s}^{\Wcond}(H)$ and the empirical fraction
of seeds with $B_{W,s}(H)\ge\beta$ over candidate thresholds.  No model was
retrained: Figure~\ref{fig:enl-learnability-profile} is derived from the tracked
per-seed E-NL summary.  Tanh terminal min/mean/max are
$4.41313266754/4.49664855003/4.67605876923$, with 8/8 above $\beta=0.5$; GRU
min/mean/max are $0.0828229486942/0.129690139554/0.187142759562$, with 0/8 above
$0.5$.  Thus the GRU weak-only gate failure is not marginal at that operating
point, but it still provides neither a causal-starvation certificate nor evidence
that starvation is absent.  Terminal reach is not generally interchangeable with
the preregistered first-hit event on a nonmonotone trajectory, and this post-hoc
diagnostic supports neither a new confirmatory endpoint nor an architecture
ranking.

\ifdefined\ICLRSubmission\else
\artifactfigure{../paper/artifacts/enl_tanh_crossover-20260824-141207/enl_learnability_profile.pdf}
{Post-hoc weak-only reach audit from the completed E-NL run.  Points are terminal
responses for all eight seeds; horizontal segments mark medians and the dashed line
marks the preregistered $\beta=0.5$.  The right panel sweeps candidate thresholds
without rerunning training.  Terminal reach is not generally interchangeable with
a first-hit gate for nonmonotone trajectories.}{fig:enl-learnability-profile}
\fi

The post-hoc tanh search over $\rho\in\{1.5,2,3,4,6\}$ and
lags $\{0,1,2,4,8\}$ found no cell with 4/4 causal certificates; the best cells
reached 3/4.  We retain this as a negative search audit and do not select a searched
cell for a paper claim.  In particular, the superseded 8/8 tanh result and its
$\tau^*=1.6112$ used a different response convention and are not valid under the
final definition.

\artifactfigure{../paper/artifacts/enl_tanh_crossover-20260824-141207/enl_crossover.pdf}
{Paired nonlinear diagnostic trajectories.  Direct-autograd drift is primary;
projected and matched-state quantities are separately styled diagnostics.}{fig:enl}

\subsection{E19 falsification; E20 restricted event-label follow-up}
The initialization-frozen tangent surrogate was first tested under a sealed five-output
pilot using one unseen data seed and eight unseen model seeds per architecture.  It
obtained 12/16 phase, 14/16 drift-crossing, 14/16 response-crossing, 14/16 causal-
certificate, and 12/16 weak-learnability classifications.  The preregistered rule
nevertheless \emph{failed}: tanh learnability was 4/8, below the 5/8 per-model
floor.  We preserve that rejection rather than averaging it away.

Only after this failure did we define the narrower crossing-event hypothesis.  A
fresh sealed factorial crossed four unseen data seeds with eight unseen model seeds
for each architecture.  Table~\ref{tab:ntk-heldout} reports classification
correctness, not observed crossing prevalence.  The restricted rule passed: all
response-crossing labels and 60/64 drift-crossing labels were correct; the two-way
bootstrap lower endpoints exceeded the frozen 0.5 floor.

\begin{table}[t]
\centering
\caption{Manifest-sealed frozen-kernel evaluations.  The broad pilot's complete
five-output rule failed.  Only the two event rows in the fresh factorial were
preregistered primary endpoints.  CIs resample data- and model-seed axes.}
\label{tab:ntk-heldout}
\small
\begin{tabular}{llcccc}
\toprule
Protocol & Endpoint & Overall & tanh & GRU & 95\% CI \\
\midrule
Broad pilot & phase & 12/16 & 6/8 & 6/8 & diagnostic \\
Broad pilot & weak learnability & 12/16 & 4/8 & 8/8 & \textbf{failed} \\
Restricted $4\times8$ & drift event & 60/64 & 32/32 & 28/32 & $[0.8125,1]$ \\
Restricted $4\times8$ & response event & 64/64 & 32/32 & 32/32 & $[1,1]$ \\
\bottomrule
\end{tabular}
\end{table}

The restriction is substantive.  In the factorial, phase, causal-certificate, and
learnability correctness were only 52/64, 56/64, and 49/64 and were explicitly
secondary; tanh learnability was 17/32.  Crossing times remained systematically
early (overall drift/response bias $-0.822/-1.189$), and tanh response-gap and weak-
response RMSE were 2.455 and 2.489 despite high sign agreement.  Thus the evidence
supports an uncalibrated event classifier at this one width/task cell.  It neither
proves Theorem~\ref{thm:frozen-comparison}'s kernel premise nor validates phase,
causal starvation, learnability, trajectory magnitude, calibrated time, or
prevalence.  Appendix~\ref{app:frozen-negative} reports the complete failure,
restricted-success, timing-bias, and magnitude-error numbers.

\subsection{Historical E3 CDC ablation: local mechanism control, not best final tradeoff}
In historical E3, all five shadow-based methods reduce the causal weak AUC gap from
ERM's 10.5155 to approximately zero and attain their shared weak-drift target at
every logged step.
Their weak outcomes are numerically similar to roughly four decimal places, but
not statistically equivalent: 9/10 direct paired comparisons are significant at
0.05, while the differences are practically tiny relative to the roughly 10.5
change from ERM.  CDC alone preserves instantaneous strong drift to rounding scale,
$4.77\times10^{-7}$, versus $1.48\times10^{-1}$ to $3.75\times10^{-1}$ for the
alternative shadow constraints.  This local property does not yield the best final
strong response: unconstrained rescue and the PCGrad-style shadow-target rescue end at
1.295, compared with CDC's 1.042 and ERM's 1.364.  These are historical E3
outcomes, distinct from E28's later frozen tradeoff.

\begin{table}[t]
\centering
\caption{Dense CDC ablation, eight paired seeds.  Shadow-based arms have access to
a matched weak-only model; the other arms do not.  ``Strong change'' is the maximum
absolute instantaneous strong-drift perturbation.}
\label{tab:cdc}
\small
\resizebox{\linewidth}{!}{%
\begin{tabular}{lcccc}
\toprule
Method & Shadow & Weak AUC gap & Strong change & Final $m_s$ \\
\midrule
ERM & no & 10.5155 & --- & 1.364 \\
Spectral Decoupling & no & 9.0229 & --- & 0.829 \\
Interaction & no & 10.5145 & --- & 1.365 \\
CDC & yes & $-0.0005$ & $4.77\times10^{-7}$ & 1.042 \\
Loss-gradient projection & yes & $-0.0002$ & $3.75\times10^{-1}$ & 0.371 \\
Unconstrained rescue & yes & $-0.0007$ & $1.48\times10^{-1}$ & 1.295 \\
Bloop-style shadow-target rescue & yes & $0.0002$ & $3.75\times10^{-1}$ & 0.370 \\
PCGrad-style shadow-target rescue & yes & $-0.0007$ & $1.48\times10^{-1}$ & 1.295 \\
\bottomrule
\end{tabular}}
\end{table}

\artifactfigure{../paper/artifacts/historical/e3_cdc_dense_ablation-20260823-090702/e3_mitigation.pdf}
{Eight-method dense-linear mitigation ablation.  CDC's established distinction is
instantaneous first-order preservation, not superior weak rescue or final strong
retention.}{fig:e3}

\subsection{Causal gate indeterminacy and comparator-restricted mechanism control}
E26 completed 64 semi-real records.  MNIST 3-vs-8 and FashionMNIST 0-vs-6 each
had 0/32 weak-only-learnable records and therefore 0/32 causal certificates under
the specified gate.  Their mean weak-AUC gaps were as follows.  For MNIST 3-vs-8,
the gap was $-1.2362842136667493\times10^{-3}$.  Its 95\% CI was
$[-3.4405428466004646,\allowbreak 1.0159591371180453]\times10^{-3}$.  For
FashionMNIST 0-vs-6, the gap was $-1.0921728159018996\times10^{-4}$.  Its 95\%
CI was $[-4.129384520886106,\allowbreak 3.182670049955049]\times10^{-3}$.
Positive means weak-only core gain exceeds both-feature core gain.  The frozen
hypothesis failed.  Because every weak-only first-hit gate failed, the causal gate
correctly rejects a false-positive starvation interpretation, but the zero causal
counts remain negative/indeterminate: they are not evidence that starvation is
absent and do not show that either dataset is intrinsically unlearnable.

E27 completed 192 nonlinear records.  The equal-weight three-cell all-$\beta$
causal-certificate macro estimate was $0.2916666666666667$ for tanh (95\% CI
$[0.16666666666666666,0.4166666666666667]$) and $0$ for GRU (95\% CI $[0,0]$);
the Holm-adjusted one-sided $p$-value was $1.0$ for each.  The overall beta-robust
claim is false.  Architecture ranking was forbidden and neither performed nor
supported.

E28 completed 256 method-records.  Define weak rescue as
$-\lvert\text{weak AUC gap}\rvert$, so the CDC-minus-comparator contrast is
higher-better.  Define trajectory and final deviation as absolute deviations from
paired ERM, so comparator-minus-CDC is positive when CDC preserves the strong
response better.  Table~\ref{tab:e28-tradeoff} reports every frozen macro contrast.
The frozen joint weak-rescue/strong-preservation tradeoff is true only versus
Bloop-style shadow-target rescue.  It is false versus PCGrad-style shadow-target
rescue, so CDC does not beat that variant under the frozen joint criterion.  Both
comparators have \code{canonical:false}; these results establish neither SOTA
performance, canonical parity, nor broad CDC superiority.

\begin{table}[t]
\centering
\caption{E28 equal-weight two-cell macro contrasts over 256 method-records.  Both
comparators are non-canonical.  Brackets are 95\% bootstrap CIs.}
\label{tab:e28-tradeoff}
\small
\resizebox{\linewidth}{!}{%
\begin{tabular}{lllll}
\toprule
Comparator & CDC$-$comparator weak rescue & Comparator$-$CDC trajectory dev. & Comparator$-$CDC final dev. & Joint \\
\midrule
Bloop-style shadow-target rescue & $0.002198259399210656$ $[0.0009779187294930126,0.003277366267916477]$ & $3.4651361294978416$ $[3.3443827215131754,3.5849715262780952]$ & $0.6422362388111651$ $[0.6288088704226539,0.6563168084365315]$ & true \\
PCGrad-style shadow-target rescue & $-0.00015351238407674823$ $[-0.00047909348592838785,0.00013248012419808214]$ & $-0.48709889128076556$ $[-0.6897378944428417,-0.27379327373499074]$ & $-0.2565436437726021$ $[-0.3016841153614223,-0.2087651835754514]$ & false \\
\bottomrule
\end{tabular}}
\end{table}

E20 remains the submission-primary empirical study, with only drift- and
response-crossing classification primary.  E26--E28 were prospectively frozen but
completed later; their outcomes do not reprioritize E20 or promote its secondary
phase, causal, learnability, timing, or magnitude outputs.

\section{Discussion and limitations}
\paragraph{Causal diagnosis and empirical findings.}
Weak learning failure alone is not gradient starvation.  E1 shows positive paired
weak-AUC deficits with learnable weak-only controls in all 32 $\rho$--seed cells
per lag at lags 0, 2, and 4, but those cells reuse eight seed IDs and do not imply a
universal law.  All 32 lag-8 cells instead fail the gate under
$\beta=0.25,H=5$, making their causal status indeterminate.  At E-NL's fixed
$(\rho,\mathrm{lag})=(4,2)$ cell, tanh has complete transfer-to-starvation
certificates in 3/8 observed seeds; GRU has drift and response crossings in 5/8 but
0/8 causal certificates because every weak-only path fails $\beta=0.5$.  The post-
hoc terminal profile shows that this GRU gate failure is not marginal at that
operating point, but it neither replaces the first-hit gate nor ranks
architectures.  These counts are descriptive, not prevalence estimates, and the
exploratory search found no 4/4 robust cell.  E26 adds semi-real breadth but is
negative/indeterminate because MNIST 3-vs-8 and FashionMNIST 0-vs-6 each have 0/32
weak-only-learnable and 0/32 causal records; the gate blocks a false-positive
starvation label but cannot establish absence or intrinsic dataset unlearnability.
E27 rejects the frozen all-$\beta$ claim for both architectures without licensing a
ranking.  E28 is mixed and comparator-restricted: its frozen joint tradeoff holds
against the non-canonical Bloop-style shadow-target rescue but not the non-
canonical PCGrad-style shadow-target rescue.  None of these later outcomes changes
E20's formal submission-primary drift/response-only scope.

\paragraph{Exact dynamics and pathwise attribution.}
The exact finite-width identity $\dot M_w=c_w^\top q(r)/n$ separates
cross-entropy sample weighting $q(r)$ from response-specific geometry $c_w$ with no
DMFT assumption.  Together with exact sample-logit dynamics, dense-linear geometry,
and the matched intervention, it identifies where a weak-response rate difference
comes from without silently imposing a two-response closure.  The cumulative
balance then gives pathwise attribution for arbitrary sign changes:
$\Delta=P-N$, so outcome suppression occurs only when accumulated suppression
overtakes accumulated help.  Neither rate nor outcome suppression is by itself a
causal-starvation certificate; independently demonstrated weak-only learnability
is still required.

\paragraph{Partial prediction and formal scope.}
The realized-initialization $d(0)$ criterion gives an exact local sign in its stated
dense-linear setting.  Where a higher local jet is proved, such as
$d(0)=0,d'(0)<0$ in the recurrence-invisible zero-lag result, it likewise gives a
rigorous local sign in that scope.  The conditional parameter-tube and rank-one
criteria, transverse hitting-time bound, recurrence-invisible zero-lag wide limit,
and fixed-chart CDC result are mathematical statements with proofs, but their
scopes differ.  The cumulative balance characterizes outcome suppression without
supplying weak-only learnability; the rank-one criterion addresses rate crossing;
the wide limit requires zero lag and zero background input, although it permits any
finite realized recurrent matrix.  Stronger positive-lag parameter-level
prediction from $\rho$, lag, and recurrent geometry remains unresolved because no
nonvacuous general product-ball or kernel-movement constants are constructed.

\paragraph{Frozen-surrogate sequence.}
The initialization-frozen tangent surrogate retains nonlinear cross-entropy but is
exact only for its tangent model.  Its comparison with trained tanh/GRU networks
assumes unproved instantaneous or secant-kernel movement control.  E19 failed its
preregistered five-endpoint rule; correctness on other labels cannot rescue the
tanh learnability failure.  E20 was designed only after that rejection and supports
a narrower claim: drift/response crossing-event classification on one fresh
crossed-seed factorial.  Its phase, causal, and learnability labels were secondary,
its time predictions were systematically early, and tanh trajectory magnitudes
were poor.  Manifest integrity proves that frozen inputs were consumed unchanged;
it does not prove kernel stability, surrogate accuracy, or general causal
prediction.

\paragraph{Local mechanism-level control and baselines.}
In the fixed implemented chart, feasible uncapped target-attaining CDC is the
unique minimum-Euclidean-norm correction satisfying zero instantaneous strong-
drift change and the weak-drift lower bound; its finite-step guarantee is local and
one-step.  Historical E3 exhibits that instantaneous distinction within its tested
shadow-method family, but does not show a general practical advantage: weak rescue
is numerically similar across shadow arms, other arms retain a larger final strong
response, and all shadow methods use stronger matched-counterfactual information.
E28 passes the separately frozen joint weak-rescue/strong-preservation criterion
only against Bloop-style shadow-target rescue; CDC does not beat PCGrad-style
shadow-target rescue under that criterion.  Both comparators are non-canonical
response-gradient shadow-target variants, not independently audited canonical
replicas.  CDC assumes a constructible weak-only shadow and roughly doubled model
work; no SOTA, canonical parity, or broad superiority follows.

\paragraph{External validity.}
The evidence covers controlled synthetic task families and one semi-real generated-
cue intervention over fixed cells and seed grids.  E26 is not an unmodified natural
shortcut or natural-spurious-correlation study.  E27/E28 broaden the fixed
synthetic cells and methods but do not establish architecture rankings, prevalence,
canonical baseline parity, or broad CDC superiority.  The unrun Waterbirds adapter
is a minibatch head-coordinate surrogate outside the full-batch matched-shadow
theorem, and its intercept is not an identified bird-shape response.  We make no
Waterbirds or natural-spurious-correlation claim.  General recurrent causal
prediction and calibrated tanh/GRU crossing times remain open.

\paragraph{Provenance.}
The final E-NL source is reconstructible through its archived 17-file executable-
source fingerprint and preserved metadata.  E19 and E20 preserve compact outcomes,
seals, and source digests, but not the exact dirty source bytes; they are auditable
but not historically replayable.  Earlier E1, E2, E2-R, exploratory E-NL, and
historical E3 share a disclosed dirty-source limitation.  E26--E28 instead ran from
clean commit \code{9ebbc617a9aaad5c3c19e9d5fd4fa08f2d3517d0}; their recorded
source manifests make source reconstruction possible.  The E26 compact archive
omits 192 \path{records/**} files and the combined E27/E28 archive omits 1,792.
Their committed hashes and central files support audit, but omitted bytes cannot be
recovered from hashes, so the compact archives are not standalone replay packages.

\section{Conclusion}
Weak learning failure alone is not gradient starvation.  A small weak response can
reflect failure of the weak-only control to meet the specified target, altered
cross-entropy weighting, response-specific geometry, rate suppression that has not
yet reversed earlier help, or causally gated outcome suppression.  The matched
same-initialization intervention separates these possibilities.  Exact finite-
width dynamics give $\dot M_w=c_w^\top q(r)/n$ without a DMFT assumption; the
pathwise identity $\Delta=P-N$ then shows that outcome suppression begins only when
accumulated suppression overtakes accumulated help.  Causal starvation additionally
requires an independently learnable weak-only counterfactual.  A failed gate
licenses neither a starvation certificate nor evidence that starvation is absent;
causal status is indeterminate under the specified target and horizon.

Prediction remains partial.  Realized $d(0)$ and proved higher local jets give
rigorous local or assumption-scoped sufficient signs.  The conditional rank-one,
parameter-tube, hitting-time, and recurrence-invisible zero-lag results state their
additional assumptions explicitly; the noiseless symmetric zero-lag reference
starts suppressed, permits arbitrary finite recurrent gain, and does not supply a
transfer phase.  Because nonvacuous general positive-lag product-ball and kernel-
movement constants are unavailable, the paper does not provide a global predictor
from $\rho$, lag, and RNN geometry.

The empirical sequence supports the diagnostic separation rather than a universal
starvation law.  Corrected E1 has positive paired weak-AUC deficits and learnable
weak-only controls across all 32 $\rho$--seed cells per lag at lags 0, 2, and 4,
not 32 independent seeds; every lag-8 cell fails the gate under
$\beta=0.25,H=5$ and is causally indeterminate.  At E-NL's fixed
$(\rho,\mathrm{lag})=(4,2)$ cell, tanh has complete transfer-to-starvation
certificates in 3/8 observed seeds, whereas GRU has drift and response crossings in
5/8 but 0/8 causal certificates because every weak-only path fails $\beta=0.5$.
The post-hoc terminal profile strengthens that operating-point diagnosis without
replacing the first-hit gate, ranking architectures, or estimating prevalence.
E19 remains a failed broad frozen-kernel pilot.  E20 remains the formal submission-primary
empirical study and supports only restricted frozen drift/response crossing-event
labels.  E26 is a completed negative/indeterminate semi-real study: MNIST 3-vs-8
and FashionMNIST 0-vs-6 each have 0/32 weak-only learnability and 0/32 causal
records, so the gate blocks a false-positive label but cannot show absent
starvation or intrinsic dataset unlearnability.  E27's all-$\beta$ claim is
negative for both architectures without a ranking.

As mechanism-level control, in the fixed implemented chart, feasible uncapped
target-attaining CDC is the unique minimum-Euclidean-norm correction satisfying
zero instantaneous strong-drift change and the weak-drift lower bound; its finite-
step guarantee is local and one-step.  Historical E3 demonstrates that local
constraint but not broad practical superiority or best final retention.  E28's
frozen joint weak-rescue/strong-preservation tradeoff succeeds only against the
non-canonical Bloop-style shadow-target rescue; CDC does not beat the non-canonical
PCGrad-style shadow-target rescue under that criterion.  Thus no SOTA,
architecture ranking, canonical parity, or broad CDC superiority follows.

The strongest supported chain is exact dynamics $\to$ matched causal intervention
$\to$ pathwise attribution $\to$ local mechanism-level control.  With fixed seed
grids, no natural-spurious-correlation benchmark, historically unreconstructible
source for E19/E20, compact E26--E28 archives that are not standalone replay
packages, and no general positive-lag DMFT, general recurrent causal prediction and
calibrated crossover times remain open.

\ifdefined\ICLRSubmission\else
\appendix
\fi
\section{Proofs for exact dynamics and causal-separation results}
\label{app:finite-proofs}

\subsection{Proof of Proposition~\ref{prop:cross-kernel} and
Corollary~\ref{cor:projected}}
For $\ell(r)=\log(1+e^{-r})$, $\ell'(r)=-\sigm(-r)$.  Therefore
\begin{equation}
 \dot\theta_C=-\nabla L_C
 =\frac1{n_C}\sum_i\sigm(-r_i^C)\nabla r_i^C.
\end{equation}
The chain rule gives
\begin{equation}
 \dot M_a^C=\ip{\nabla M_a^C}{\dot\theta_C}
 =\frac1{n_C}\sum_i\ip{\nabla M_a^C}{\nabla r_i^C}\sigm(-r_i^C),
\end{equation}
which proves Proposition~\ref{prop:cross-kernel}.  Under exact realization,
$\nabla r_i^C=\sum_bz_{ib}^C\nabla M_b^C$.  Substitution and exchange of finite
sums yield Eq.~\eqref{eq:Gg}.  If an orthogonal residual logit term is present,
substitution of
$\nabla r_i^C=\sum_bz_{ib}^C\nabla M_b^C+\nabla r_i^\perp$ instead yields
Eq.~\eqref{eq:projection-residual}.  \qed

\subsection{Proof of Corollary~\ref{cor:full-empirical-ntk} and
Theorem~\ref{thm:frozen-comparison}}
The parameter flow above can be written
$\dot\theta=J^\top q(r)/n$.  Therefore the chain rule gives
$\dot r=J\dot\theta=JJ^\top q(r)/n$ and
$\dot M=\ip{\nabla M}{\dot\theta}=c^\top q(r)/n$, proving
Corollary~\ref{cor:full-empirical-ntk}.

For the comparison, $q$ is $1/4$-Lipschitz componentwise and hence in Euclidean
norm.  Subtracting Eq.~\eqref{eq:frozen-surrogate} from
Eq.~\eqref{eq:full-ntk-flow}, adding and subtracting
$\Theta_0q(r)$, and writing $e_r=\norm{r-\bar r}$ gives
\begin{equation}
 e_r'(u)\le \frac{\norm{\Theta_0}}{4n}e_r(u)+\frac{B_\Theta V}{n}
 =ae_r(u)+\frac{B_\Theta V}{n}.
\end{equation}
Gr\"onwall with $e_r(0)=0$ proves Eq.~\eqref{eq:frozen-logit-bound}.  The same
add--subtract step gives
\begin{equation}
 |\dot M-\dot{\bar M}|\le \frac{B_cV}{n}
 +\frac{\norm{c_0}}{4n}e_r.
\end{equation}
Integrating and using $\int_0^t\phi_a(u)\dd u=\psi_a(t)$ proves
Eq.~\eqref{eq:frozen-response-bound}.

For SGD, the fundamental theorem of calculus along the straight update segment
gives
\begin{align}
 r_{j+1}-r_j&=\widetilde J_j\delta\theta_j
 =\frac{\eta}{n}\widetilde\Theta_jq(r_j),\\
 M_{j+1}-M_j&=\ip{\widetilde m_j}{\delta\theta_j}
 =\frac{\eta}{n}\widetilde c_j^\top q(r_j).
\end{align}
Subtracting the frozen Euler recursion yields scalar recurrences
\begin{align}
 e_{j+1}&\le(1+\eta a)e_j+\frac{\eta D_\Theta V}{n},\\
 E_{j+1}&\le E_j+\frac{\eta D_cV}{n}
 +\frac{\eta\norm{c_0}}{4n}e_j.
\end{align}
Solving the first geometric recurrence and summing it in the second gives
Eq.~\eqref{eq:frozen-discrete-bound}, including the stated $a=0$ limits.  Paired
bounds are the triangle inequality.  Error-separated sign transfer follows
immediately; first-hit transfer requires the additional assumptions of
Corollary~\ref{cor:hitting}.  \qed

\subsection{Proof of the dense-linear geometry}
For a unit response mode in channel $a$ at lag $\ell_a$, write $b_a$ for the
corresponding input column.  Direct differentiation of
$c^\top W^{\ell_a}b_a$ gives gradient $p_a=W^{\ell_a}b_a$ with respect to $c$,
gradient $q_a=(W^{\ell_a})^\top c$ with respect to column $b_a$, and recurrent
gradient
\begin{equation}
 S_a=\sum_{j=0}^{\ell_a-1}(W^j)^\top c
       (W^{\ell_a-1-j}b_a)^\top.
\end{equation}
The parameter blocks are disjoint, so their gradient inner products add.  Input
columns contribute only when the channels coincide, giving
Eq.~\eqref{eq:dense-geometry}.  \qed

\subsection{Proof of Theorem~\ref{thm:initial-gap}}
For the positive rank-one cue laws,
$g^{\Bcond}=s_B(\rho,1)$ and $g^{\Wcond}=s_W(0,1)$ at the shared
initialization.  Corollary~\ref{cor:projected} and the exact dense geometry give
\begin{equation}
 \dot M_w^{\Bcond}(0)=s_B(\rho G_{ws}+G_{ww}),\qquad
 \dot M_w^{\Wcond}(0)=s_WG_{ww}.
\end{equation}
Subtraction, with $K=G_{ws}$ and $A=G_{ww}$, proves
Eq.~\eqref{eq:initial-gap-certificate}.  Since $\rho s_B>0$, rearrangement gives
Eq.~\eqref{eq:initial-gap-threshold}.  The displayed $K,A,S_w$ formulas are
Eq.~\eqref{eq:dense-geometry} specialized to a zero-lag strong mode and an
$\ell$-lag weak mode.  Under the iid width scaling at fixed positive $\ell$,
expanding the gates at zero gives
\begin{equation}
 d_0=\rho\left(\frac K2-\frac{AM_s(0)}4\right)+O_{\Prob}(N^{-1}),
\end{equation}
whose leading sign is centered and realization-dependent.  Shared initialization
also gives $\Delta(0)=0$ and $\Delta'(0)=d_0$.  Therefore
$\Delta(t)/t\to d_0$ as $t\downarrow0$: a strict negative (positive) $d_0$
forces strict local outcome suppression (transfer) on some sufficiently small
positive interval.  This limit establishes no quantitative endpoint, and when
$d_0=0$ first order is silent.  This proves the exact criterion, its strongest
unconditional local outcome implication, and the stated obstruction to a
distribution-only positive-lag boundary.  \qed

\subsection{Proof of Proposition~\ref{prop:drift-balance} and
Corollaries~\ref{cor:tail-area}--\ref{cor:crossing-bounds}}
The positive and negative parts satisfy $d=d_+-d_-$ almost everywhere.
Absolute continuity and the fundamental theorem of calculus therefore give
\begin{equation}
 \Delta(t)=\Delta(0)+\int_0^t d(u)\dd u=P(t)-N(t),
\end{equation}
proving Proposition~\ref{prop:drift-balance} and its pointwise sign equivalences.
Under Corollary~\ref{cor:tail-area}'s sign pattern, $d_-=0$ before $\tau_d$
and $d_+=0$ afterward.  Hence
$P(t)=P(\tau_d)=A_+$ and
\begin{equation}
 \Delta(t)=A_+-N(t),\qquad t\ge\tau_d.
\end{equation}
The post-crossing sign assumption makes $N$ continuous and strictly increasing on
every nontrivial post-crossing interval.  Later equality, strict negativity, and
uniqueness are therefore respectively equivalent to $N(H)\ge A_+$,
$N(H)>A_+$, and strict post-crossing growth of $N$.

If $d\le-\kappa$ after $t_0$, then
$N(t)-N(t_0)\ge\kappa(t-t_0)$; solve for the time needed to cover the remaining
area.  If $d'(t)\le-\lambda$ and $d(\tau_d)=0$, integration gives
$d(t)\le-\lambda(t-\tau_d)$ and then
$N(t)-N(\tau_d)\ge\lambda(t-\tau_d)^2/2$.  Since
$N(\tau_d)=0$ under the single-crossing hypotheses, applying
Corollary~\ref{cor:tail-area} proves both stated bounds.  \qed

\subsection{Proof of Theorem~\ref{thm:parameter-tube} and
Corollary~\ref{cor:parameter-safe}}
Before exit from $U_R$, path length is at most $S_Rt$, so exit cannot occur before
$h_R=R/S_R$.  Along this interval, $d'=j$, and integration of
Eq.~\eqref{eq:parameter-tube-premises} gives
\begin{equation}
 d_0-\Lambda t\le d(t)\le d_0-\lambda t.
 \label{eq:parameter-tube-linear-bounds}
\end{equation}
The lower bound is positive before $d_0/\Lambda$, the upper bound is negative
after $d_0/\lambda$, continuity gives a zero between them, and $d'<0$ makes it
unique.  This proves Eq.~\eqref{eq:parameter-tube-crossing}.  Integrating the upper
bound from $\Delta(0)=0$ gives
\begin{equation}
 \Delta(t)\le d_0t-\frac{\lambda t^2}{2}<0
 \quad\text{for }t>2d_0/\lambda.
\end{equation}
For the required nondegenerate case $\beta>M_w^{\Wcond}(0)$, integrating the
weak-only lower drift bound reaches $\beta$ by the positive time in
Eq.~\eqref{eq:parameter-tube-learnability}, proving the causal conclusion.  If
$\beta\le M_w^{\Wcond}(0)$, attainment occurred before optimization and supplies
no learnability gate.

For the safe case, $|j|\le J$ gives $d(t)\ge d_0-Jt$; one more integration gives
Eq.~\eqref{eq:parameter-safe}.  The strict premise $d_0-Jh_R>0$ makes both
quantities positive on the certified horizon.  \qed

\subsection{Proof of Theorem~\ref{thm:log-ratio} and
Corollary~\ref{cor:unique-rate-crossing}}
The deterministic cue law gives
$g^{\Bcond}=s_B(\rho,1)$ and $g^{\Wcond}=s_W(0,1)$.  Equation~\eqref{eq:Gg}
therefore gives $F_w^{\Bcond}=s_BH_B$, $F_w^{\Wcond}=s_WH_W$,
$u_B'=s_BQ_B$, and $u_W'=s_WH_W$.  Since
\begin{equation}
 \frac{\dd}{\dd u}\log\sigm(-u)=-(1-\sigm(-u)),
\end{equation}
differentiating $\Psi$ yields Eq.~\eqref{eq:Psi-prime}.  When both drifts are
positive, monotonicity of the logarithm gives
$\operatorname{sign}(F_w^{\Bcond}-F_w^{\Wcond})=\operatorname{sign}(\Psi)$.

Under $\Psi'\le-\kappa$, integration gives
$\Psi(t)\le\Psi(t_0)-\kappa(t-t_0)$.  Continuity forces a zero by
$t_0+\Psi(t_0)/\kappa\le H$, while strict decrease gives uniqueness.  The
primitive inequalities imply the derivative bound term by term.  \qed

\subsection{Proof of Corollary~\ref{cor:hitting}}
Continuity and first-entry semantics imply $x(\tau)=\beta$.  Absolute continuity
and $x'\ge\kappa$ almost everywhere give, by integration,
\begin{equation}
 x(t)\le\beta-\kappa(\tau-t)\quad(\tau-r\le t\le\tau),
 \qquad
 x(t)\ge\beta+\kappa(t-\tau)\quad(\tau\le t\le\tau+r).
 \label{eq:hitting-cones}
\end{equation}
On $[0,\tau-r]$, the quantitative prehistory margin and $\epsilon<\eta$
give $x_N<\beta$.  On
$[\tau-r,\tau-\epsilon/\kappa)$, the left cone and the uniform error again give
$x_N<\beta$.  Hence $\tau_N\ge\tau-\epsilon/\kappa$.  At
$t=\tau+\epsilon/\kappa$, the right cone gives $x_N(t)\ge\beta$, so the hit is
finite and $\tau_N\le\tau+\epsilon/\kappa$.  This proves
Eq.~\eqref{eq:hitting-continuous}.

For sampled data, the first grid point at or after
$\tau+\epsilon/\kappa$ is at most $h$ later and, by the radius condition, remains
inside the right cone; it is therefore at or above $\beta$.  No sampled hit can
precede the continuous lower bound.  If the piecewise-linear interpolant first
crosses on $[t_{j-1},t_j]$, then $x_N(t_j)\ge\beta$, so
$t_j\ge\tau-\epsilon/\kappa$ and
$t_{j-1}\ge\tau-\epsilon/\kappa-h$; its upper bound is no later than the sampled
hit.  This proves Eq.~\eqref{eq:hitting-grid}.  Applying the bound separately to
two finite hits and using the triangle inequality proves the paired statement.
\qed

\subsection{Proof of Theorems~\ref{thm:six-scalar} and
\ref{thm:wide-limit}}
At zero lag and zero background input, every cue occurs at the final step.  From
$h_0=0$, induction gives zero hidden state before that step, so for every finite
realized $W$,
\begin{equation}
 h_T=Wh_{T-1}+Bx_{T-1}=Bx_{T-1}.
\end{equation}
The loss is independent of $W$, each recurrent-gradient contribution contains a
zero incoming state, and therefore $\nabla_WL=0$ throughout the unregularized
flow.  No distributional or zero-gain premise on $W$ is needed.  Gradient flow for
the remaining parameters gives $\dot b_a=g_ac$ and
$\dot c=\sum_bg_bb_b$.  Differentiating $M_a=c\cdot b_a$,
$u_{ab}=b_a\cdot b_b$, and $w=c\cdot c$ proves Eq.~\eqref{eq:six-scalar}.
The mode gradient has one component in the $c$ block and one in its own $b_a$
block, yielding $G_{ab}=u_{ab}+\delta_{ab}w$.

At initialization, the variances of $M_s,M_w,u_{sw}$ are $1/N$ and those of
$u_{ss}-1,u_{ww}-1,w-1$ are $2/N$.  Coordinatewise Chebyshev bounds and a union
bound sum to $1+1+1+2+2+2=9$, proving
Eq.~\eqref{eq:init-concentration}.

The field bound $\norm{g(M)}_\infty\le A$ follows from
$0\le\sigm\le1$.  Since $\sigm$ is $1/4$-Lipschitz,
\begin{equation}
 |g_a(M)-g_a(\widetilde M)|
 \le\frac14\sum_b\E|z_az_b|\,\norm{M-\widetilde M}_\infty,
\end{equation}
which proves Eq.~\eqref{eq:g-lipschitz}.  Define the Gram energy
$R=w+u_{ss}+u_{ww}$.  Along the deterministic flow,
\begin{equation}
 R'=4g\cdot M=4\E[X\sigm(-X)],\qquad X=z\cdot M.
\end{equation}
Because $x\sigm(-x)\le c_\star=W(e^{-1})$ and $R(0)=3$,
$R(t)\le R_H$.  Every coordinate of $q$ is a squared norm or inner product of
$b_s,b_w,c$, so Cauchy--Schwarz gives $\norm{q(t)}_\infty\le R(t)\le R_H$.

Consider two states in the cube $\norm q_\infty,\norm{\widetilde q}_\infty
\le B_{H,r}$ and set $d=\norm{q-\widetilde q}_\infty$.  Every product in the
$M$ and $u$ equations satisfies the direct estimate
\begin{equation}
 |g_a(M)y-g_a(\widetilde M)\widetilde y|
 \le (A+B_{H,r}C)d.
 \label{eq:product-lipschitz}
\end{equation}
There are at most three such products in an $M$ equation and two in a $u$
equation; the factor two and two cue indices give four in the $w$ equation.
Thus the full vector field is $L_{H,r}=4(A+B_{H,r}C)$-Lipschitz on the cube and
therefore on the radius-$r$ reference tube.

Let $T_r$ be the first time before $H$ that
$\norm{q_N(t)-q(t)}_\infty$ reaches the tube boundary.  Up to $T_r$, Gronwall's
inequality and Eq.~\eqref{eq:tube-bootstrap-event} give
\begin{equation}
 \norm{q_N(t)-q(t)}_\infty
 \le e^{L_{H,r}t}\norm{q_N(0)-q_0}_\infty.
\end{equation}
If $L_{H,r}>0$, the initialization event makes this bound strictly less than $r$
for every $t<H$, so the boundary cannot be reached before $H$.  If $L_{H,r}=0$,
then $A=C=0$, the cue law is zero almost surely, and the six-scalar vector field
vanishes; the initial discrepancy is constant and cannot exit the closed tube.
The endpoint may lie on the boundary.  In either case the same estimate is at most
$\delta$ throughout $[0,H]$, proving Eq.~\eqref{eq:gronwall}.  Applying
Eq.~\eqref{eq:init-concentration} with
$\epsilon=e^{-L_{H,r}H}\min\{r,\delta\}$ proves
Eq.~\eqref{eq:trajectory-concentration}.  For the shared-initialization pair, use
a common $r$ and the maximum deterministic constant from the two fixed cue laws;
the same concentrating initial six scalars control both paths.  \qed

\subsection{Proof of Corollary~\ref{cor:zero-disorder-suppression}}
Under the stated noiseless positive rank-one cue law, the recurrence-invisible
six-scalar equations give equal initial weak drifts in the two conditions for every
finite realized $W$.  Differentiating those drifts once at the symmetric state gives
\begin{equation}
 (\dot M_w^{\Bcond})'(0)=-\frac{\rho^2+1}{2},\qquad
 (\dot M_w^{\Wcond})'(0)=-\frac12,
\end{equation}
so $\bar d'(0)=-\rho^2/2$.  Taylor's theorem and continuity then make
$\bar\Delta(t)<0$ for all sufficiently small positive $t$.

The weak-only symmetric subsystem preserves $u_{ww}=w$.  Direct differentiation
shows $u_{ww}^2-M_w^2=1$, yielding
$u_{ww}=w=\sqrt{1+M_w^2}$ and the displayed positive drift.  While
$0\le M_w\le\beta$, it obeys $\dot M_w\ge2\sigm(-\beta)$; integration proves
Eq.~\eqref{eq:zero-disorder-target-time}.  The finite-width statement is then
exactly Theorem~\ref{thm:wide-limit} applied at deterministic suppression and
target-attainment witness times with margins smaller than both strict witness
margins.  \qed

\section{Proofs for local Counterfactual Drift Correction}
\label{app:cdc-proofs}
\subsection{Instantaneous preservation}
If $a=0$, both strong drifts vanish.  Otherwise,
\begin{equation}
 a\cdot q=a\cdot p-\frac{p\cdot a}{\norm a^2}(a\cdot a)=a\cdot p-p\cdot a=0,
\end{equation}
so $a\cdot v'=a\cdot v+\alpha a\cdot q=a\cdot v$.  This argument is independent
of $\alpha$, hence survives a cap.  \qed

\subsection{Unique minimum-norm correction}
Let $\bar\delta=F_w^{\Wcond}-F_w$ be the raw deficit and
$\delta=[\bar\delta]_+$.  If $\bar\delta\le0$, $d=0$ is feasible and uniquely
norm-minimal.  Suppose $\bar\delta>0$, so $\delta=\bar\delta$.  Any feasible $d$
obeys $a\cdot d=0$ and $p\cdot d\ge\delta$.  Because
$p=q+(p\cdot a/\norm a^2)a$ when $a\ne0$ (and $q=p$ when $a=0$),
$p\cdot d=q\cdot d$.  Cauchy--Schwarz in the fixed Euclidean subspace orthogonal
to $a$ gives
\begin{equation}
 \delta\le q\cdot d\le\norm q\norm d,
 \qquad \norm d\ge\delta/\norm q.
\end{equation}
Equality requires a nonnegative multiple of $q$.
The correction $d^*=(\delta/\norm q^2)q$ is feasible and attains this lower
bound, so it is the unique minimum-norm, lower-bound-attaining feasible correction.
Larger positive multiples can also attain the target but have larger norm.  A
positive deficit is feasible exactly when $q\ne0$ in this representation.  A
nonbinding cap leaves $d^*$ unchanged; a binding cap generally loses target
attainment, so the target-attaining optimization claim no longer applies.  \qed

\subsection{Local finite-step deviation}
In the fixed chart and norm, $L$-Lipschitzness of the gradient on both straight
segments gives the smoothness remainder
\begin{equation}
 |m_s(\theta+h)-m_s(\theta)-\nabla m_s(\theta)\cdot h|
 \le\frac L2\norm h^2.
\end{equation}
Apply this at $h=\eta v$ and $h=\eta v'$.  Instantaneous preservation cancels the
linear terms and the triangle inequality yields Eq.~\eqref{eq:cdc-step} for the
actual velocities, capped or uncapped.  For active uncapped CDC,
\begin{equation}
 \norm{v'}\le\norm v+\frac{\delta}{\norm q}\le V+D/q_{\min},
\end{equation}
which gives the explicit bound; a nonnegative upper cap can only decrease the
added multiple of $q$.  If smoothness is known only on the fixed-coordinate ball
$B(\theta,r)$, sufficient segment conditions are
$\eta V\le r$ and $\eta(V+D/q_{\min})\le r$.  These quantities and $L$ are
chart- and metric-dependent.  \qed

\section{Discrete certificate and experiment audit}
\label{app:experiment-audit}
For logged values $\Delta_k$, increments telescope exactly:
\begin{equation}
 \Delta_K-\Delta_0=\sum_{k=0}^{K-1}(\Delta_{k+1}-\Delta_k).
\end{equation}
The executable certificate requires one strict block of positive increments followed
by one strict block of negative increments.  The peak-minus-initial value is its
positive area; peak-minus-final is its negative tail area.  It certifies causal
starvation only if the final trajectory is strictly below the initial gap and the
external weak-only learnability flag is true.  A false learnability flag blocks the
causal certificate but does not establish that starvation is absent; status is
indeterminate under that gate.  This finite-step certificate makes no assertion
between logged points.

\paragraph{Excluded engineering outcomes.}
Data seed 100 with model seeds 1000--1007 is excluded from scientific evidence.  A
first frozen-prediction grid mismatch stopped before training; a second observed-
trajectory grid mismatch was detected after nonlinear outcomes had been generated,
but no scores were inspected.  The set was used only to repair evaluator alignment
through integer-step matching plus a numerical optimization-time check.  It did not
change labels, thresholds, or the five E19 endpoints; E19 and E20 use fresh seeds.

\begin{table}[ht]
\centering
\caption{Executable contracts in \code{src/gradient\_starvation/theory.py}.}
\label{tab:contracts}
\small
\resizebox{\linewidth}{!}{%
\begin{tabular}{lll}
\toprule
Contract & Mathematical object & Scope guard \\
\midrule
\code{projected\_statistics} & direct Eq.~\eqref{eq:cross-kernel}, $Gg$, residual & direct drift primary for nonlinear probes \\
\code{initial\_frozen\_empirical\_kernel} & $J_0,\Theta_0,c_0$ in Corollary~\ref{cor:full-empirical-ntk} & initialization only; full sample kernel \\
\code{integrate\_frozen\_logistic\_sgd} & explicit Euler for Eq.~\eqref{eq:frozen-surrogate} & exact for tangent surrogate, not nonlinear model \\
\code{frozen\_kernel\_discrete\_error\_bound} & Eq.~\eqref{eq:frozen-discrete-bound} & caller supplies secant-kernel movement bounds \\
\code{equal\_time\_drift\_difference} & $\Delta'$ and two projected product orderings & exposes non-canonical attribution \\
\code{dense\_linear\_initial\_gap\_certificate} & Eqs.~\eqref{eq:initial-gap-certificate}--\eqref{eq:initial-gap-threshold} & noiseless positive rank-one dense-linear scope only \\
\code{parameter\_tube\_certificate} & Theorem~\ref{thm:parameter-tube}, Corollary~\ref{cor:parameter-safe} & caller supplies uniform product-ball constants; missing premises fail closed \\
\code{rank\_one\_drift\_ratio} & Theorem~\ref{thm:log-ratio} factors and $\Psi$ & caller supplies noiseless/exact hypotheses \\
\code{discrete\_crossover\_certificate} & telescoping finite-step tail areas & learnability gate required for causal label \\
\code{transverse\_hitting\_time\_error\_bound} & $\epsilon/\kappa+h$ & supplied error, prehistory, radius, and mesh guards \\
\code{zero\_disorder\_*\_failure\_bound} & Eqs.~\eqref{eq:init-concentration}, \eqref{eq:trajectory-concentration} & caller supplies reference-tube constants \\
\code{cdc\_finite\_step\_deviation\_bound} & Eq.~\eqref{eq:cdc-step} & caller supplies local smoothness and segment containment \\
\bottomrule
\end{tabular}}
\end{table}

\begin{table}[ht]
\centering
\caption{Executed experiment configurations.  ``Shadow'' methods train a second
weak-only model and therefore use roughly twice the model memory and per-step
compute.}
\label{tab:protocol}
\small
\resizebox{\linewidth}{!}{%
\begin{tabular}{lllllll}
\toprule
Experiment & Model / width & Task grid & Samples & Steps / LR & Seeds / records & Primary purpose \\
\midrule
E1 & dense linear / 96 & $\rho\in\{1,2,4,8\}$, lag $\in\{0,2,4,8\}$ & 256 & 500 / .01 & 4--11 & paired deficit / gate audit \\
E2 & dense linear / 32--256 & $(1,0),(2,2),(4,4),(8,8)$ & 1024 & 400 / .01 & 0--7 & empirical reference \\
E-NL & tanh, GRU / 32 & $\rho=4$, lag 2 & 256 & 1000 / .01 & 20--27 & direct-drift certificate \\
Tangent pilot & tanh, GRU / 32 & $\rho=4$, lag 2 & 256 & 1000 / .01 & $1\times8$ & five-output falsification \\
Tangent factorial (E20) & tanh, GRU / 32 & $\rho=4$, lag 2 & 256 & 1000 / .01 & $4\times8$ / model & primary crossing-event audit \\
Historical E3 & dense linear / 96 & $\rho=4$, lag 2 & 256 & 500 / .01 & 40--47 & CDC local ablation \\
E26 & semi-real CNN & MNIST 3-vs-8; FashionMNIST 0-vs-6 & 512/256/256 per dataset & 500 / .01 & $4\times8$ / dataset; 64 & generated-cue intervention \\
E27 & tanh, GRU / 32 & three cells; $\beta\in\{.25,.5,1,2\}$ & 256 & 1000 / .01 & $4\times8$; 192 & all-$\beta$ robustness \\
E28 & tanh / 64 & two cells; ERM, CDC, Bloop-style shadow-target rescue, PCGrad-style shadow-target rescue & 256 & 1000 / .01 & $4\times8$; 256 method-records & comparator-restricted CDC tradeoff \\
\bottomrule
\end{tabular}}
\end{table}

\ifdefined\ICLRSubmission
\paragraph{Supplementary artifact map and commands.}
The anonymous supplementary archive separates the controlled paired-deficit and
gate audit,
finite-width reference, exact solver checks, final nonlinear run and terminal
$\beta$ audit, failed broad tangent pilot, restricted E20 tangent factorial,
exploratory negative search, historical E3 CDC ablation, E26 semi-real
negative/indeterminate study, E27 beta-robust negative study, and E28 mixed CDC
tradeoff.  E20 remains submission-primary with only drift/response crossings
primary.  Machine-readable provenance binds each claim to configurations, outcomes,
seals, and source status.  Run \code{python verify\_release\_artifacts.py} for the
read-only compact-archive audit.  New scientific runs
use \code{scripts/run\_semi\_real\_study.py} or
\code{scripts/run\_expanded\_studies.py} only after fresh preflight and external
seal; they are not replays of omitted record bytes.  E26--E28 use reconstructible
clean source at commit \code{9ebbc617...2d3517d0}, but the E26 and combined E27/E28
compact archives omit 192 and 1,792 \path{records/**} files, respectively, so they
are auditable rather than standalone replay packages.  E19/E20 instead lack their
historical dirty source bytes and are not historically replayable.
\else
\begin{table}[ht]
\centering
\caption{Authoritative artifact map.  Paths are relative to the repository root.}
\label{tab:artifacts}
\small
\begin{tabular}{@{}p{0.38\linewidth}p{0.57\linewidth}@{}}
\toprule
Claim family & Artifact directory \\
\midrule
Controlled paired deficits and gate status & \path{results/e1_dense_rerun-20260823-085142} \\
Finite-width empirical closure & \path{results/e2_corrected_validation-20260821-224630} \\
Exact special solver checks & \path{results/e2r_solver_checks-20260823-084911} \\
Final nonlinear crossover and $\beta$ audit & \path{paper/artifacts/enl_tanh_crossover-20260824-141207} \\
Broad frozen-NTK pilot (failed) & \path{paper/artifacts/enl_ntk_pilot-20260825} \\
Restricted frozen-NTK factorial (crossing labels passed) & \path{paper/artifacts/enl_ntk_crossing_factorial-20260825} \\
Exploratory nonlinear audit & \path{results/enl_exact_regime_search-20260824-105003} \\
Historical dense CDC ablation & \path{results/e3_cdc_dense_ablation-20260823-090702} \\
E26 semi-real generated cue (negative/indeterminate) & \path{paper/artifacts/semi-real-generated-cue-v1-20260830} \\
E27 beta-robust nonlinear (negative) & \path{paper/artifacts/expanded-studies-v1-20260830} \\
E28 CDC tradeoff (comparator-restricted mixed) & \path{paper/artifacts/expanded-studies-v1-20260830} \\
\bottomrule
\end{tabular}
\end{table}

\paragraph{Commands for new runs under the current source.}
The following commands execute new runs from the repository root; they do not
bitwise reproduce historical dirty-source outcomes:
\begin{verbatim}
python run_experiment.py e1  --config configs/e1_dense_rerun.yaml
python run_experiment.py e2  --config configs/e2_corrected_validation.yaml
python run_experiment.py e2r --config configs/e2r_solver_checks.yaml
python run_experiment.py enl --config configs/enl_tanh_crossover.yaml
python scripts/build_enl_learnability_profile.py
python run_experiment.py e3  --config configs/e3_cdc_dense_ablation.yaml
python verify_release_artifacts.py
\end{verbatim}
The last command is a read-only audit of the committed compact archives.  Fresh
E26-style execution uses
\begin{verbatim}
python scripts/run_semi_real_study.py preflight \
  --data-root DATA --output-dir PREFLIGHT
\end{verbatim}
followed, only with an external seal, by
\begin{verbatim}
python scripts/run_semi_real_study.py run \
  --preflight PREFLIGHT --seal SEAL \
  --data-root DATA --output-dir RUN
\end{verbatim}
Fresh E27/E28-style execution similarly uses
\begin{verbatim}
python scripts/run_expanded_studies.py preflight \
  --output-dir PREFLIGHT
\end{verbatim}
and the sealed \code{run} subcommand.  Such runs are new studies, not replays of
the compact archives' omitted 192 or 1,792 \path{records/**} files.
The checked-in frozen evaluate configs are historical evidence records: they pin the
original preflights and intentionally fail against later source.  A new protocol run
must first execute \code{enl-preflight}, then copy the matching evaluate config and
explicitly repin \code{evaluation.preflight\_dir} and
\code{evaluation.manifest\_sha256} before \code{enl-evaluate}.  This creates a new
study; it does not bitwise reproduce the archived dirty-source studies.  The broad
pilot pins manifest SHA-256
\code{f186d587...f0ffaf5} and prediction SHA-256
\code{9787d999...f9dfaa}; its evaluation verdict is failed.  The restricted
factorial pins manifest SHA-256 \code{bdcf02ab...a0db1} and prediction SHA-256
\code{c05b3181...6542};
its crossing-only verdict is passed.  The exploratory E-NL command is deliberately
omitted from the confirmatory list; its complete configuration and negative result
remain archived in the repository.  The learnability-profile command is a
read-only derivation from the tracked E-NL summary and does not run training.
\fi

\section{Completed causal-gate and mechanism-control outcome and provenance audit}
\label{app:completed-extensions}
These studies were prospectively frozen, externally authorized, and executed from
clean Git commit \code{9ebbc617a9aaad5c3c19e9d5fd4fa08f2d3517d0}.  They are
later breadth studies, not an outcome-dependent replacement for E20: E20 remains
the submission-primary empirical claim and only its drift- and response-crossing
classification endpoints are primary.

\subsection{E26: semi-real causal indeterminacy after weak-only gate failure}
E26 contains 64 crossed records: four data seeds by eight model seeds for MNIST
3-vs-8 and the same factorial for FashionMNIST 0-vs-6.  The core channel contains
selected real benchmark pixels; the paired intervention exactly zeros only the
generated cue channel.  Thus the evidence label is semi-real generated-cue, not an
unmodified real-world or natural-spurious-correlation benchmark.  Positive weak
AUC gap means the weak-only core gain exceeds the both-feature core gain.

\begin{table}[ht]
\centering
\caption{Complete E26 primary outcomes.  CIs are dataset-specific crossed-seed
intervals.}
\label{tab:e26-complete}
\small
\resizebox{\linewidth}{!}{%
\begin{tabular}{lccc}
\toprule
Dataset & Weak-only learnable & Causal certificate & Mean weak AUC gap (95\% CI) \\
\midrule
MNIST 3-vs-8 & 0/32 & 0/32 & $-0.0012362842136667493$
$[-0.0034405428466004646,0.0010159591371180453]$ \\
FashionMNIST 0-vs-6 & 0/32 & 0/32 & $-0.00010921728159018996$
$[-0.004129384520886106,0.003182670049955049]$ \\
\bottomrule
\end{tabular}}
\end{table}

The frozen hypothesis was not accepted.  Because the weak-only first-hit gate
failed in all 32 records per dataset, the causal gate rejects a false-positive
starvation interpretation.  The zero causal counts remain negative/indeterminate:
they cannot be interpreted as evidence that starvation is absent or that either
dataset is intrinsically unlearnable.  The compact archive appears under
\path{paper/artifacts/} as \path{semi-real-generated-cue-v1-20260830}; the central
provenance manifest records its exact archive, artifact, and 22-file clean-source
digests.
Source is reconstructible from the clean commit, but the compact archive omits 192
\path{records/**} files; hashes bind those files but cannot recover their bytes, so
the archive is not a standalone replay package.

\subsection{E27: expanded beta-robust nonlinear study}
E27 contains 192 records: tanh and GRU across three fixed task cells and a crossed
$4\times8$ seed design.  Each record evaluates the prospectively frozen first-hit
grid $\beta\in\{0.25,0.5,1,2\}$; the architecture estimand equally weights the
three cell prevalences.  A 5,000-replicate two-way pigeonhole bootstrap supplies
raw percentile intervals, and the two one-sided architecture tests form one Holm
family.

\begin{table}[ht]
\centering
\caption{Complete E27 all-$\beta$ architecture endpoints.}
\label{tab:e27-complete}
\small
\resizebox{\linewidth}{!}{%
\begin{tabular}{lccc}
\toprule
Architecture & Macro estimate & 95\% CI & Holm-adjusted $p$ \\
\midrule
tanh & $0.2916666666666667$ & $[0.16666666666666666,0.4166666666666667]$ & $1.0$ \\
GRU & $0$ & $[0,0]$ & $1.0$ \\
\bottomrule
\end{tabular}}
\end{table}

The overall beta-robust claim is false.  Architecture ranking was forbidden by the
frozen contract and was neither performed nor supported.

\subsection{E28: comparator-restricted mechanism-level control}
E28 contains 256 method-records for ERM, CDC, Bloop-style shadow-target rescue,
and PCGrad-style shadow-target rescue over two fixed cells and crossed seeds.  Both
comparators are non-canonical (\code{canonical:false}) response-gradient
shadow-target variants; the study does not establish canonical implementation
parity.  Weak rescue is $-\lvert\text{weak AUC gap}\rvert$ and is higher-better, so
CDC-minus-comparator positive values favor CDC.  Trajectory deviation is absolute
strong-response AUC deviation from paired ERM and final deviation is absolute final
strong-response deviation from paired ERM; both are lower-better, so
comparator-minus-CDC positive values favor CDC.

\begin{table}[ht]
\centering
\caption{Complete E28 frozen equal-weight two-cell macro contrasts.  Brackets are
95\% two-way-bootstrap CIs.}
\label{tab:e28-complete}
\small
\resizebox{\linewidth}{!}{%
\begin{tabular}{lllll}
\toprule
Comparator & CDC$-$comparator weak rescue & Comparator$-$CDC trajectory deviation & Comparator$-$CDC final deviation & Joint tradeoff \\
\midrule
Bloop-style shadow-target rescue & $0.002198259399210656$ $[0.0009779187294930126,0.003277366267916477]$ & $3.4651361294978416$ $[3.3443827215131754,3.5849715262780952]$ & $0.6422362388111651$ $[0.6288088704226539,0.6563168084365315]$ & true \\
PCGrad-style shadow-target rescue & $-0.00015351238407674823$ $[-0.00047909348592838785,0.00013248012419808214]$ & $-0.48709889128076556$ $[-0.6897378944428417,-0.27379327373499074]$ & $-0.2565436437726021$ $[-0.3016841153614223,-0.2087651835754514]$ & false \\
\bottomrule
\end{tabular}}
\end{table}

The frozen joint weak-rescue/strong-preservation tradeoff is true only against
Bloop-style shadow-target rescue and false against PCGrad-style shadow-target
rescue.  CDC therefore does not beat the latter under the frozen joint criterion.
Because both comparators are non-canonical, E28 supports no SOTA, canonical parity,
or broad CDC superiority.  E27 and E28 share
\path{paper/artifacts/expanded-studies-v1-20260830}; the central provenance manifest
records its exact archive, artifact, and 23-file clean-source digests.  Source is
reconstructible from the clean commit, but the compact archive omits 1,792
\path{records/**} files and is not a standalone replay package.

\section{Frozen-tangent falsification and restricted follow-up}
\label{app:frozen-negative}
The sealed E19 pilot tested one preregistered conjunction over five outputs and 16
architecture--seed records.  Correct classifications were phase 12/16, drift
crossing 14/16, response crossing 14/16, causal certificate 14/16, and weak-only
learnability 12/16.  The conjunction failed because tanh learnability was 4/8,
below the frozen 5/8 floor; GRU learnability was 8/8.  Magnitude calibration also
differed sharply by architecture: tanh response-gap, weak-response, and drift RMSE
were 2.367, 2.394, and 0.486, versus 0.0163, 0.0281, and 0.0107 for GRU.  These
numbers are preserved as a falsification result rather than averaged into a pass.

E20 was specified only after E19's rejection and used fresh crossed seeds.  Its
accepted scope was restricted to event labels: response crossing was correct in
64/64 records and drift crossing in 60/64 (32/32 tanh, 28/32 GRU), with two-way
bootstrap intervals $[1,1]$ and $[0.8125,1]$.  The rejected secondary outputs were
substantially weaker: phase 52/64, causal certificate 56/64, and learnability 49/64,
including tanh learnability 17/32.  Predictions were systematically early
(overall drift/response bias $-0.822/-1.189$), while tanh response-gap,
weak-response, and drift RMSE were 2.455, 2.489, and 0.490.  Consequently E20 is an
uncalibrated crossing-event classification result at one controlled cell, not a
phase, causal, learnability, time, magnitude, prevalence, or kernel-stability
validation.  Both compact archives preserve these scores and seals but not the
unavailable exact dirty source snapshots, so integrity auditability does not imply
historical replayability.

\ifdefined\ICLRSubmission\else
\bibliographystyle{plainnat}
\bibliography{references}
\fi

\end{document}
